concilio, y que así puede pecar o errar; pues aunque como hombres sean falibles, ha de creerse que el verdadero concilio legítimamente reunido no se rige por sentido humano, sino por espíritu divino, de donde están suficientemente probadas las cosas que el concilio aprueba, y no han de ser discutidas ulteriormente por la singularidad y testarudez de cualquiera, sino que debemos cautivar el entendimiento en obediencia de la fe. Por lo cual me maravillo de que, estando por negar el texto de los derechos, admita como comentarista al Panormitano.

Así, en cuanto a lo que por último adujo de las palabras del apóstol Pablo sobre el fundamento, siempre he dicho que es sentir de todos los católicos que Cristo es el fundamento de la Iglesia, que Cristo es la cabeza de la Iglesia, pues él es el Señor principal. Pero porque constituyó a Pedro como vicario y le confirió potestad sobre toda la Iglesia, como atestigua san León en el can. Ita Dominus 19. dist.. Y me maravillo sobremanera de que, dejado a un lado el texto de san León, el reverendo padre recoja allí una glosa contraria al texto. Convenía en todo caso, al ascender Cristo a los cielos, que en su lugar constituyese aquí un vicario como cabeza, al que se recurriese en las dudas y en los demás casos emergentes, para que toda la Iglesia de Cristo no fuese un «uno» confuso, un caos anaxagórico. Esta sentencia la sostiene toda la Iglesia, como está en el responsorio con que la Iglesia habla de san Pedro: Tú eres el pastor de las ovejas, príncipe de los apóstoles. A ti te entregó Dios todos los reinos del mundo. Así oramos en la Iglesia. Más aún, el divino Francisco, aprobado con los sellos de los estigmas, enseñó a sus hermanos y a todo el orbe de la tierra a obedecer al sumo pontífice, y esto lo preceptuó en la regla. Por eso canta de él la Iglesia: Francisco, varón católico y todo apostólico, enseñó a guardar la fe de la Rom. Iglesia. Así muchísimos sumos pontífices, así los sagrados concilios lo definieron, cosas que por brevedad no aduzco. Muchísimas sentencias de los santos padres suenan y aprueban lo mismo. Todo lo cual no descuidaré aducir a su tiempo, cuando vea el tratado del reverendo padre sobre este asunto.

Si, sin embargo, el pri-\marginpar{\small[297]}mado del Rom. pontífice se apoya sólo en derecho humano y es elevado por consenso del pueblo, ¿de dónde entonces tiene el reverendo padre el privilegio de la mendicidad? ¿De dónde ese hábito de la religión? ¿De dónde la potestad

de oír confesiones, de predicar, con otros innumerables privilegios, «con los que son abrumados por la sede apostólica», contradiciéndolo con frecuencia arzobispos, obispos y curas? Los cuales, sin embargo, como hijos de la obediencia escuchan al Rom. pontífice y en él a Cristo, para alcanzar así por el bien de la paciencia la vida eterna.

Estas breves cosas he querido añadir en el presente asunto de la disputa, según me lo dio el tiempo, habiendo de aducir por lo demás más, si no me arredrara el fastidio y el mandato del ilustrísimo príncipe.

\paragraph*{El D. D. Martín.}

Dos cosas me incumben hacer: primero responder a las confutaciones del egregio D. D., segundo objetar contra él según el pacto de la convención. Y porque ya durante tres días he respondido y se ha consumido el tiempo, de modo que ni una sola puedo cumplir, digo brevemente:

Que todo lo aducido por el egregio D. D. ha sido dicho ciertamente con grandísima copiosidad, pero del todo nada al propósito, pues durante tres días prometió que actuaría por derecho divino contra mí, y hasta ahora no se ha oído ni una autoridad de derecho divino, fuera de aquella de Mt 16, que todos oímos, de cómo diversamente es tratada por los santos padres, y de que la mayor y más sana parte de ellos sostiene la sentencia en mi favor. Él mismo adujo en su favor un solo lugar de Agustín y otro de Ambrosio; todo lo demás fueron palabras de Anacleto, Marcelo, Pelagio etc., de quienes escribieron derechos humanos. Cipriano, en cambio, casi todo siente conmigo, lo cual encomiendo al juicio del prudente lector. Igualmente «el clarísimo texto» de Jerónimo a Evagrio, y sobre Tito. Luego Gregorio en el registro por casi seis epístolas seguidas, después de esto la historia y el uso de la iglesia oriental. Pues nada hace al caso que el D. D. haya alegado a los griegos cismáticos, «a Nestorio, Macedonio, Acacio, Eutiques y demás», pues aquéllos no eran la iglesia oriental. Podría «también» yo describir así a la iglesia latina por medio de «pelagianos, maniqueos, jovinianos», vigilantianos y monstruos semejantes. Pero líbreme de llamar cismática a toda una iglesia por unos pocos malos y cismáticos.

Lo demás «no podré proseguirlo». Y me reservo especialmente los artículos de Juan Hus a mí objetados «y entre ellos el novísimo sobre los actos humanos» para el cálamo y el papel, y todo lo demás que tenía

que objetar, para prestar por escrito lo que aquí no puedo hacer, excluido por los límites de espacio, y quiero en esto haber cedido y obedecido a los mandatos del ilustrísimo príncipe, clementísimo patrono de los estudios, de lo cual os requiero a vosotros, señores notarios, y llamo por testigos a los oyentes.

Hecho como arriba, presentes allí los honorables Olavo Roters de Osvaberga y Celestino Goritio de Gorlitz, clérigos de las diócesis de Colonia y de Meissen, testigos especialmente llamados y legítimamente requeridos.

Juan Graumann, notario, lo suscribió.

\marginpar{\small[298]} El jueves, día siete del mes de julio, a la hora séptima de la mañana.

\paragraph*{El D. D. Eck.}

Porque el reverendo padre tuvo ayer la hora de responder a nuestras cosas, para que así hubiésemos ejecutado los mandatos del ilustrísimo príncipe, donde al responder apeló muchísimo por la palabra «pase» al cálamo, sin embargo, habida deliberación, propuso que ya respondería más plenamente. Yo no me opuse, por amor a la verdad y para su esclarecimiento. Me admiro de cómo el reverendo padre con gran estruendo se atreve a tronar ante tantos varones que yo habría dicho que probaría por derecho divino la monarquía y el primado en la Iglesia, cuando sin embargo sólo alegué dichos de los santos padres y decretos de los pontífices, cuando sin embargo yo le aduje el derecho divino de Mt 16 sobre la promesa de la edificación sobre la piedra, de que Cristo rogó especialmente por la fe de Pedro, de que especialmente le encomendó confirmar a los hermanos, de que especialmente le encomendó en presencia de Juan y de Santiago apacentar las ovejas, por tres veces. Y que este derecho divino ha de entenderse así, lo probé entonces con los santos padres, los santos sumos pontífices y mártires, el consenso de los concilios y de todas las universidades. Por lo cual por derecho divino quedó probado que hay primado en la Iglesia, a no ser que el reverendo padre anteponga su propia inteligencia al entendimiento de tantos santos padres, sumos pontífices y del sacro concilio, pues el laudable concilio Constanciense condenó tal artículo entre los pestilentes errores de Juan Hus. A lo cual el reverendo padre calla. Pero sin embargo, si quiere ofrecer cosas mejores, para esclarecimiento de la verdad, lo que ayer habría podido

haber hecho según el mandato del ilustrísimo príncipe, estoy dispuesto a oírlo, con el pacto de que, mientras desempeña el oficio de respondiente, ya al fin del acto no se traslade taimadamente a la persona del oponente, sino que cualquier cosa que haya de oponer contra esta ineludible verdad, estoy dispuesto a responder a estas cosas aducidas y a disolverlas, de lo cual protesto.

\paragraph*{El D. D. Martín.}

Como ayer se había establecido en nombre del ilustrísimo príncipe, nuestro patrono, que esta materia se consumase aquel día, y el egregio D. D. hubiese consumido sin necesidad y más odiosamente de lo que convenía a este espléndido auditorio la mayor parte del tiempo, a mí para una respuesta de tres días sólo me quedó una horilla, en la que a la vez respondiese a tanta farrago suya y objetase. Plugó después por gratuito afecto al mismo ilustrísimo príncipe nuestro que se me diese a mí copia de responder y de objetar, por lo cual doy a su ilustrísima gracia las máximas gracias. Y por eso, viniendo al propósito, antes de continuar la comenzada respuesta de ayer, disiparé primero esta niebla matutina del egregio D. D., donde a su modo con bastante odiosidad siempre me objeta las pravidades husitas, no ob-\marginpar{\small[299]}stante que tantas veces me he purgado de ellas. Y por eso prologaré un poco en alemán, porque entiendo que soy muy mal oído entre el vulgo.

Decía pues que no impugnaba, ni podía ser cristianamente reprobado por nadie el primado y la obediencia de la Romana iglesia. Sino que había sido llevado al punto de tener que declarar que ese primado no es de derecho divino, así como tampoco puede reprobarse la potestad imperial entre los germanos, aunque no fundada en la sagrada escritura.

\paragraph*{El D. D. Eck.}

dijo que era verdad que ésa era la controversia, si aquel primado es de derecho divino, cosa que juzgaba reprobar como husítica y artículo condenado por el concilio Constanciense, y que habría que dolerse si los corazones de los cristianos fuesen tan fríos que nadie defendiese esto.

Después el D. D. Martín prosiguió en latín diciendo.

Que en cuanto aduce esta autoridad de Mt 16, contendiendo que allí fue prometido el primado, bastante se ha respondido en lo precedente, que

el egregio D. apenas tiene dos autores y esos mismos ambiguos para su sentencia, cuando sin embargo la mayor parte establece que Pedro fue allí la persona de todos los apóstoles y de los fieles, cosa que también urge el mismo tenor consecuente del texto, pues Cristo interroga a todos, y no sólo a Pedro, diciendo: Mas vosotros ¿quién decís que soy yo? Y allí Pedro, boca de los apóstoles, como dice Crisóstomo, y persona de los apóstoles, como dicen Jerónimo y Orígenes y Beda y Agustín y Ambrosio, responde por todos. Por eso dije que nada se prueba de esta autoridad sobre el primado. Igualmente, que el texto mismo no dice así: Tú eres Pedro, y sobre ti, sino sobre esta piedra, indicando manifiestamente otra piedra, de la que Pedro fue dicho. Por lo demás, como son derechos divulgados y los mismos Rom. pontífices atestiguan que la interpretación de la escritura vale más en los doctores que en ellos mismos, pero en las causas por decidir al contrario.

Sobre aquel artículo husítico dije bastante ayer, que aún no se ha probado que sea herético. Esto por respuesta matutina.

Volviendo a lo de ayer, cuando el egregio D. D. en segundo lugar opuso que la deposición de Crisóstomo fue hecha ilegítimamente porque se hizo por odio de la emperatriz, niego la consecuencia; no fue por eso ilegítima la potestad de ordenar y deponer establecida por el sacro concilio Niceno, porque por acaso tuvo mal uso. Pues en la misma historia se escribe que el Rom. pontífice Víctor mártir quiso excomulgar a los obispos de Asia. Pero ellos a su vez como superiores les mandaron silencio, y les ordenaron no perturbar a la iglesia de Dios, lib. 5 de la historia eclesiástica. Pero también Ireneo de Lyon de la Galia, obispo, coerció al mismo Rom. pontífice, y eran entonces tem-\marginpar{\small[300]}porales en Grecia Epifanio, celebradísimamente alabado por Jerónimo, Gregorio Nacianceno, Basilio, cristianísimos varones, y sin embargo nunca estuvieron bajo el Rom. pontífice, sino que según los estatutos del concilio Niceno eran ordenados por los obispos de las provincias.

En cuanto a que el egregio D. doctor aduce que Julio I excomulgó al emperador Arcadio, nada hace al caso. Pues también Bonifacio VIII se atrevió a querer arrojar de la sede al rey de los francos. No prueba que sea derecho todo lo que los Rom. pontífices por humana fragilidad alguna vez intentaron. Ni es maravilla que aquellos santos varones fuesen vejados por la tentación de la ambición

cuando más de una vez los santísimos apóstoles aun en presencia de Cristo padecieron lo mismo.

En cuanto a que dice que al obispo «jerosolimitano» le fue dada prerrogativa de honor, no primado, rectamente lo dice, pero no prueba que haya primado del Rom. pontífice, al menos por derecho divino.

También allí dice óptimamente el egregio D. D. que el sexto sínodo dio a la Romana iglesia el primado, pero no los sínodos precedentes. Esto es pues lo que yo quise, que por estatutos sinodales y por derecho humano, con el consenso de los fieles de Cristo, al que no es lícito resistir, fue dado ese primado. Como también el divino Gregorio dice en el registro con estas palabras: Sabéis que el venerable sínodo Calcedonense ofreció a esta sede apostólica el primado y sin embargo ninguno de mis predecesores presumió jamás aceptarlo. Por lo demás nada hace que el egregio D. D. diga que Gregorio no rechazó el primado, sino que lo afirmó. Lea el piadoso lector las epístolas de Gregorio, y hallará que Gregorio prohibió el primado y no vindicó el pontificado ecuménico, como arriba por lo relatado en la dist. 99. bastante probamos.

A la autoridad de Agustín, a Inocencio Rom. pontífice y a la de éste a aquél respondo: Lea el lector las palabras de uno y otro y juzgue si el egregio D. D. impuso nuevas significaciones a los vocablos.

Sobre los veinte años de la iglesia antes de la Romana dijo, habiendo sido trasladado con Pedro el primado a Roma, queriendo quizá con este proverbio: Donde está el papa, allí está Roma. Respondo: había que probar el primado que se trasladaría, no la traslación, que de buen grado concedo.

Adujo también la autoridad de Agustín contra los maniqueos, que sospecho es del libro 2 de la doctrina cristiana c. 8, que en cosa dudosa para certeza de la fe vale la autoridad de la iglesia católica, la cual es firmada por las sucesiones de los obispos que desde las mismas sedes de los apóstoles se suceden entre sí y por el consenso del pueblo. Nada dice de la Rom. iglesia, sino de la universal, y hace por mí contra el D. D., porque Agustín describe las sucesiones de las diversas sedes apostólicas entre sí y el consenso del pueblo, pero no la confirmación de las otras sedes por una sola sede.

Al asunto mismo.

\marginpar{\small[301]} Sobre la autoridad de Mt 16: Tú eres Pedro, y sobre esta piedra etc. y a los autores aducidos dice que no respondí, cosa que encomiendo

a vosotros oyentes y lectores, si es verdad. Lo de Agustín en las retractaciones, dejando la elección al lector sobre dos sentencias suyas de la misma autoridad, lo opuso diciendo que abraza ambas sentencias, lo cual me place, y así nada prueban contra mí.

La segunda autoridad de derecho divino aducida por su dignación es ésta: Apacienta mis ovejas, Jn último [21 v. 15 ff.], y me gozo de oír por fin tras tres días alguna vez una autoridad de derecho divino, en que firmísimamente se apoye su sentencia, pero veamos también ella misma.

Primero esta autoridad parece entenderse de dos modos, de un modo por el verbo amar, esto que es, fiarse de sí mismo, y que todo le sea lícito, y por el verbo apacentar, lo mismo que ser el primero y reinar. Y así será el sentido: Pedro, si me amas, esto es, si buscaras lo tuyo e hicieres todo lo que place a tus aduladores, apacienta mis ovejas, esto es, sé el primero y señor de todos. Este sentido no lo tengo en mi códice.

De otro modo: Si me amas, esto es, si te niegas a ti mismo, si pones tu alma por mí, si aun toda dignidad desdeñares y no amares sino a mí, como egregiamente en este lugar expone Agustín: Apacienta mis ovejas, esto es, enseña, predica la palabra, exhorta, ora, precede con buen ejemplo. Pues la palabra griega en este lugar no significa simplemente regir y apacentar, sino cuidar suave y blandamente y hacer todo para que nada falte a las ovejas, y creo que este sentido es el evangélico.

Por lo demás ruego al D. D. Eck que ore conmigo al Señor Jesucristo, para que quiera inspirar no sólo al sumo, sino a todos los pontífices, que crean que estas palabras les atañen. No hay duda de que todo el orbe recibiría con manos abiertas y lágrimas derramadas a tal varón que quisiese regirse según estas palabras. Por lo demás, si el egregio D. D. inspecciona diligentemente el texto, no se concede allí derecho ni privilegio, como se piensa, sino que se impone precepto, y se injunge un trabajo hoy intolerable, y se encomienda un oficio a quien de hecho está constituido ya en el primero ya en el medio lugar, y no se constituye primero. Que si urge que el oficio no puede administrarse si primero no se tiene el lugar de la potestad, place, pero entonces con otra palabra ha de construirse ese lugar, pues aquí esta palabra no encomienda sino el oficio.

Por lo demás rectamente hicieron los santos padres, que dicen que este lugar atañe a todos

pues esto nadie puede negarlo, si a Pedro le fueron aquí encomendadas todas las ovejas, porque él no apacentó a todas, como dice el bienaventurado Agustín y el libro de los Hch evidentemente demuestra, y así fue inobediente a Cristo. Ni puede decirse que, aunque por sí no apacentó, sin embargo lo hizo por otros subordinados. Primero, por conceder lo que podemos devolver, no fue dicho a Pedro: Apacienta por medio de otro, sino a él mismo en propia persona: Apacienta tú, y sin embargo esto convence de que ni ningún apóstol ni al menos Pablo fueron subordinados a Pedro, quienes apacentaron muchas más ovejas \marginpar{\small[302]} que Pedro. Por eso la inteligencia del verbo ha de pedirse de toda la escritura y de la circunstancia de los hechos, porque esta autoridad o no atañe sólo a Pedro o no a todos.

En tercer lugar digo, si por estas palabras se probara el primado, será incierto y nulo el primado, porque no se pone el primado sino bajo condición de amor, pues no conviene romper la escritura en diversas partes, sino más bien concertarla. Y por eso, como es incierto quién ama a Cristo, también nos será incierto quién sea nuestro pastor. Y si a él no le es necesario el amor, tampoco nos será necesario a nosotros reconocer al pastor. Más rectamente pues se dice que esta palabra es un precepto dado a todos los prelados de la iglesia, para que, despreciadas las riquezas, la dignidad, aun el mismo primado, por fin con vida y muerte se gasten por las ovejas de Cristo. ¿Quién es éste, y lo alabaremos?

Paso lo que adujo de Crisóstomo sobre el vértice del colegio, y cómo a Pedro fue entregada la presidencia del orbe de la tierra, porque todo lo concedo, que sea el primero en honor, pero no solo en la administración, al menos por derecho divino.

Sobre el caso de Pedro, Gál 2, donde de nuevo el egregio D. D. más odiosamente me objetó el artículo husítico de que ningún señor espiritual o civil que esté en pecado mortal —si dijese que Pedro por eso no es prelado porque pecó mortalmente—, respondo: Esto no se trataba, y esta disputa está fuera del punto [ad rhombum]. Sé que el prelado malo no ha de ser rechazado. Por eso también yo condeno aquel artículo husítico. Pero esto quise, que Pedro, cuando en causa de fe dio escándalo, si no hubiese sido enmendado por Pablo, con mérito debió ser removido de la prelatura, pues el pastor hereje o el que simula en peligro

de la fe, peca gravemente. Pues por esta simulación de Pedro perecía desde el fundamento la fe de Cristo, como dice Pablo, por eso antes habría de hacerse, si Pedro no hubiese sido enmendado, que Pedro fuese tenido por nada, mucho menos por sumo pontífice. El pecado mortal daña a la propia persona, pero la herejía daña a la persona común y a la comunidad. Por eso doy gracias al señor D. de que al menos esto aprendí de esta disputa, que edificar sobre la piedra significa estar bajo el prelado, sea bueno, sea malo. Pues esta significación de la edificación y de la piedra confieso que antes ignoraba, porque dijo: aun si Pedro hubiese caído mortalmente, aún habría permanecido piedra, esto es primero.

En cuanto a que dice que antaño los obispos no eran confirmados si no prestada cédula de fe, lo admito, estoy contento, pero nada al derecho divino, pues nada derogo al pontífice, cualquier cosa que se le atribuya.

Después de esto adujo a Nicolás Cusano y excusó a Anacleto, varón santo, porque no le fue necesaria tanta ciencia a tan santo varón. Lo confieso. Pero al sumo pontífice, sea santo, sea no santo, le es torpe ignorar los evangelios, siendo pastor de las ovejas y doctor del evangelio. Como Cefas es interpretado por Juan evangelista cap. 1 [v. 42] diciendo Pedro: Tú eres Simón, tú \marginpar{\small[303]} serás llamado Cefas, lo que se interpreta Pedro, cuando entonces el hijo del trueno así truena, el que conoce todas las lenguas, con mérito debe callar la vanidad del terrenal cordel, sea Nicolás Cusano, caldeo o hebreo. No porque reprenda al egregio D. D., que dijo que Pedro vale lo mismo que cabeza de casa en caldeo, por autoridad de su Cusano, sino porque alabo al impertérrito disputador, que osa meterse en la arena teológica y ni siquiera estudió antes tanto en los evangelios que leyese el primer capítulo del evangelio de Juan. Igualmente no tiene en su libro qué es Cefas, lo cual perdono, pero lea el capítulo primero de 1 Cor [v. 12] y lo hallará.

Igualmente en último lugar sobre los artículos bohemios no quiere juzgar, aunque nunca desiste de denigrarme, y yo sobre ellos paso, sin decir nada, sino que a Gregorio de Rímini en la dist. XXVIII, reprobado por el egregio D. D., yo lo apruebo. Pues es todo él nada más que Agustín y la divina escritura, resistiendo ciertamente a todos los doctores escolásticos, entonces máximamente a Aristóteles, pero aún por ninguno confutado.

Convengo con el señor D. en que los estatutos de los concilios en estas cosas que son de fe han de abrazarse de todo modo, sólo esto me reservo, que también ha de guardarse que el concilio alguna vez erró y puede errar, sobre todo en estas cosas que no son de fe, ni tiene el concilio autoridad de forjar nuevos artículos en la fe, de otro modo tendremos al fin tantos artículos cuantas opiniones de hombres.

En cuanto a que al fin adujo a san Francisco, que debemos obedecer al sumo pontífice, me maravillo contra quién lo dijo. Igualmente paso estas criminaciones sobre los hermanos mendicantes, este miserable hábito también atormenta al D. doctor. Yo por mi sentencia pronuncio que quiero que no haya ninguna orden de mendicantes. Esto tuve que responder a lo objetado por el egregio D. D. sobre la proposición 13. Por lo cual, acabada ahora la respuesta, resta que también yo mueva contra el egregio D. D. las cosas que me mueven.

Aquí el señor doctor Martín quiso pasar a objetar por las partes de oponente a él devueltas por pacto, como decía. Pero el señor doctor Eck pretendía que entre lo soluble ya se habían dicho algunas cosas a modo de oponente, las cuales primero querría disolver según la protestación sobre esto hecha. No obstante permitió que proceda ya la oposición del señor doctor Martín. Comenzó pues

\paragraph*{El D. D. Martín.}

Contra la sentencia del egregio D. D. parece militar lo que con bastante fuerza saltó al confutar, y se escribe en 1 Cor 3 [v. 5], donde Pablo, rebajando \marginpar{\small[304]} la persona de todos los apóstoles, dice: ¿Qué es Apolo? ¿qué Pablo? ¿qué Cefas? ministros por quienes creísteis. Donde Pablo quiere del todo quitar la ocasión de contención y de cisma; niega que ninguna iglesia pertenezca a ningún apóstol, sino que todo es común a todos. De donde al fin del tercer capítulo [3. v. 21 ff.] dice: Todo es vuestro, sea Cefas, sea Apolo, sea Pablo, mas vosotros de Cristo.

No vale la evasión con que se dice que contendieron no sobre la principalidad, sino sobre la singularidad de las personas, siendo manifiesto el texto que contendieron sobre las dignidades de las personas, pues unos preferían a Pedro sobre todos, otros a Pablo, otros a Apolo, otros en cambio solo a Cristo. Se confirma por aquello de Gal 1 y 2, donde los Gálatas, seducidos por semejante contienda a causa de la primacía de Pedro que se les encomiaba, abandonaron a Pablo y a su doctrina como si fuese más indigna. Contra lo cual prueba con largo texto que nada importa al caso la majestad de Pedro o de los otros apóstoles, diciendo que él no fue enviado a ellos ni por hombre ni mediante hombre, ni siquiera vio a Pedro ni aprendió de él, sino que todo lo tuvo y lo transmitió sin Pedro. Si, pues, la autoridad de Pedro fue necesaria como derecho divino, Pablo será en este lugar manifiestamente impío y blasfemo, como quien no quiere que se vea que ha sido enviado ni siquiera por Dios mediante hombre, y rechaza por completo la autoridad de Pedro. En tercer lugar, y más abajo con mayor claridad, donde dice: A mí los que parecían ser algo nada me comunicaron, y cómo hayan sido alguna vez nada me importa, pues Dios no acepta la persona del hombre. He aquí que dice claramente que la cualidad de Pedro y de los otros apóstoles nada le importa, lo cual sería impíísimo decir si la cualidad de Pedro hubiese de guardarse por derecho divino. Con la misma impiedad diría: Dios no acepta la persona del hombre, cuando el derecho divino y lo que es de derecho divino lo preceptúa incluso bajo pena de eterna maldición. Por lo cual parece que esta primacía y majestad, o con cualquier nombre que se designe la persona o cualidad de Pedro, no se sostiene por derecho divino.

El D. D. Eck quiso responder ya repitiendo al aire las objeciones. Entonces, por mandato de los señores transmitido por el alguacil público, fue interrumpida la disputa y suspendida hasta la hora segunda de la tarde.

Hora segunda de la tarde del mismo día.

\paragraph*{El D. D. Eck.}

Cuando el reverendo padre aduce contra mí tres autoridades del bienaventurado Pablo, ventiladas tiempo ha por los adversarios de la sede apostólica y aclaradas por los defensores de la sede apostólica, fácilmente le respondo, añadiendo los malos afeites con que se esfuerza por paliar su dogma y untar a los hombres la boca y los ojos. Digo lo primero, que no pasé por alto, como me imputa, el lugar del Apóstol 1 Cor 3 [v. 5.], aunque él mismo lo adujo falsamente con Pablo y Apolo a Cefas, porque el texto desde el principio no lo contiene, sino que al final hace mención de Cefas, 'sino

al punto' la glosa interlineal confiesa que él es mayor que todos. No reciba, pues, el reverendo padre la Sagrada Escritura por su propia cabeza, sino según los santos.

Y digo que mi solución dada fue óptima e irrefutable, y nada vale el afeite del reverendo padre, con que se burla de los sencillos; testigo es el bienaventurado Jerónimo en el libro 1 contra Joviniano, col. 18 para mí, donde dice expresamente que sobre Pedro fue edificada la Iglesia. Allí pasó por alto con fuerza el reverendo padre y nada dijo. Y más abajo: para que, constituida la cabeza, se quitase \marginpar{\small[305]} la ocasión del cisma. Luego lo que el apóstol Pablo quiso evitar —los cismas y las contiendas—, con ello no quitó el primado de la cabeza, pues los cismas 'no pueden quitarse mejor' sino por la cabeza. De donde ni una palabra hace ni un ápice toca al propósito sobre la negación del primado. Por lo cual falsamente ha comenzado hoy diciendo que apenas aduje dos s. padres y ellos ambiguos, cuando aduje padres ciertos e indudables: a Agustín, a Jerónimo, a Crisóstomo, a s. León, a Bernardo etc.. Y en cuanto a que me alaba por abrazar ambas sentencias de Agustín en el libro 1 de las retractaciones cap. 21, y afirma que él mismo hace. . . . .

El D. D. Martín niega haberlo hecho etc., salvo por desliz de lengua y de modo poco premeditado etc..

Y en cuanto a lo que en segundo lugar aduce a Pablo a los Gálatas, como si aquél hubiese escrito contra Pedro, queriendo confirmar a los Gálatas en la fe, y mostrase en todo aquel proceso que él fue igual a Pedro y a los demás, porque no lo recibió ni por hombre ni mediante hombre etc., digo que, si el reverendo padre fuese tan diligente inquisidor de la verdad como promete, habría aprendido los verdaderos fundamentos y la inteligencia de las Escrituras, de modo que no sería preciso que tan jactanciosamente dijese en medio de tan grande asamblea que él solo quiere resistir a mil. Así están la verdad y el asunto. En el apostolado y en la dignidad del orden todos los apóstoles fueron iguales; esto quiso la Sagrada Escritura, esto quisieron los santos padres. A Pablo dice Jerónimo que no habría reprendido a Pedro si se supiera desigual. Así Anacleto en el cap. In novo dijo que los apóstoles fueron dotados de igual potestad. Así Cipriano, así también los demás. Pues cuando

Pablo fue designado apóstol por Cristo, ingenuamente dijo que no lo recibió ni por hombre ni mediante hombre, y que Santiago y Pedro nada le comunicaron, pues todo esto ha de referirse al apostolado.

Otra cosa, en cambio, es la potestad de régimen y de administración, donde es concorde la sentencia de los santos padres y del sagrado concilio contra Juan Hus, que tal primado perteneció al mismo Pedro. No usa, pues, rectamente los testimonios de la Sagrada Escritura, sino que los arrastra a otra parte a su arbitrio, contra lo que el Espíritu Santo exige. Pues lo que hoy ponderaba —que Pablo dijo que nada le importaba cómo hubiesen sido— no significa esto que quería el reverendo padre, sino que significaba que Pedro y Santiago y los demás fueron idiotas y poco peritos antes de andar con el Señor, y así Dios no acepta personas.

Pero dirá, como es mordazuelo, que esto es un invento eckiano. Lea a Ambrosio en el mismo lugar, quien sin corromper la sentencia de Pablo atestigua ingenuamente que la Iglesia fue fundada sobre Pedro. Por lo cual, si no adujeré cosas más firmes, ruego que desista del intento. De donde hoy, cuando yo había dicho que es artículo husita que la obediencia eclesiástica no esté expresada en la autoridad de la Sagrada Escritura, respondió con brevedad que el artículo no es herético. Preferiría tratar punto por punto 'con el reverendo padre' y preguntar: Este es un artículo condenado, y él mismo lo sostiene, luego sostiene un artículo condenado y reprobado.

Asimismo le pregunto: si el artículo no es herético, entonces o es temerario o sedicioso u ofensivo de los oídos piadosos; comoquiera que \marginpar{\small[306]} lo conceda el mismo padre, será o temerario o 'sedicioso' o sospechoso en la herejía, u ofende los oídos piadosos contra la 'institución' teológica.

Maravilla con qué color pavonino adorna las palabras de Cristo en Mt 16. Que Cristo interrogó a todos los discípulos y que así Pedro respondió por todos, ¿quién lo niega? Pero de su propia harina añadió que por eso a Pedro 'por sí y' por los demás apóstoles le fue dicho: Tú eres Pedro, y sobre esta piedra edificaré etc., lo cual, como invento luterano, no recibo, pues se opone a los santos padres.

Sobre Crisóstomo, maravilla qué injuria hace al santo padre, como si hubiese sido depuesto legítimamente y según los decretos del concilio Niceno, siendo cosa averiguada por las historias que después de Eusebio 'de Antioquía', cismático, fueron expulsados con otros rebeldes al Romano pontífice los obispos cristianos que, buscando refugio en la primera sede, padecieron repulsa de parte de los impíos cismáticos, hasta los 'tiempos de Teodosio, emperador cristiano. Cuando los cismáticos hubieron ocupado la sede treinta años, san Gregorio Nacianceno recibió en Constantinopla el 'episcopado', no 'por ambición', sino para que no le sucediera otro cismático.

Al 'santo' Crisóstomo, expulsado injustamente por medio de Epifanio, sustituido por el cismático Arsacio, lo restituyó la autoridad del Romano pontífice. Y no vale lo que entonces calumnia el reverendo padre contra el Romano pontífice y contra los santos que se le adhirieron, como si hubiesen procedido contra aquellos cismáticos de hecho, no de derecho, cosa que dice con sus solas palabras desnudas y que tantos santos, aun encomiados por el señor padre, no consintieron con el Romano pontífice; de donde sin causa empieza ya el reverendo padre a acusar de ambición a aquellos santos padres que vivieron hace mil años, vicio del que durante tantos siglos han sido tenidos por inmunes por toda la Iglesia.

Así es capcioso y se burla de los sencillos aquello de que el bienaventurado Gregorio 'no quiso llamarse ecuménico' y pastor universal, pues ayer tanto apreció a Platina, siguiendo, 'como decía', la verdad de la historia. Ya rechaza a Platina a su arbitrio, por juzgar que en todas partes le es lícito.

Alega el decreto en su favor; si se le alega, lo rechaza, y teniendo contra sí textos abiertísimos aun de León, cap. Ita dominus 19. dist., con otros semejantes, acepta la glosa, como 'los pretendientes de Penélope', por decirlo en broma, que no pudiendo poseer a la señora seguían a las esclavas. Cierto es que también los santos pontífices escribieron llamándose obispos universales, como Sixto, como Víctor, y como en el concilio Calcedonense fue aclamado a León: Viva el santísimo León, patriarca ecuménico, aunque los pontífices, siguiendo la forma de la humildad y de Cristo, se abstuvieron de aquel nombre, confesándose mejor siervos de los siervos, y para que no se creyese que querían suprimir la autoridad de los obispos, como tiempo ha 'respondió' a aquello santo Tomás en el libro sobre la impugnación de la fe.

Sobre la carta de Agustín os plantea a vosotros, señores auditores, un embrollo, para que leáis las palabras; las palabras oídas fueron: Inocencio dijo hace mil cien años que es antigua costumbre consultar al sumo pontífice en aquello que es de la fe \marginpar{\small[307]}. Esto dicho contra los maniqueos quiso el reverendo padre trasplantarlo a su huerto. Pero no lo hará, si lee la carta de Agustín, en la cual no enumera la serie de la universidad de los obispos, sino la de los Romanos pontífices desde el mismo Pedro hasta sus tiempos.

Dice que la solución dada por mí sobre los veinte años no está probada por mí; agudo es el disputador 'reverendo padre; esto aún no sabía, que el respondiente debe probar su solución; siempre juzgué que el respondiente debe ofrecer la solución para que el oponente la impugne. Débil fue su Hércules', si se apoya en Marte ajeno.

Pero que con tanta impudencia se atreva a sugerir a vuestras ilustrísimas dominaciones y reverendas paternidades que yo no leí el primer capítulo de Juan, es insigne temeridad en verdad, cuando siendo niño aún no decenal, quitados los profetas, leí toda la Biblia; pero esto nada importa al caso, cuánto haya leído cada uno.

Dije defendiendo a Anacleto que no 'fue tan ignorante aquella interpretación' de Cefas, esto es, cabeza, no ignorando que κεφαλη en griego significa cabeza en latín, sino porque en Juan, a quien ayer vi, en Agustín y en la glosa ordinaria Cefas se interpreta Pedro, y Pedro, con el Cusano por autor, también significa cabeza de la casa; todo lo cual arroja a risa el reverendo padre, sobre lo cual juzguen los jueces que han de ordenar.

Pero aquello es más excelente, que la autoridad «Apacienta mis ovejas» la envuelve en tantos rodeos, solo para hacer que los hombres viendo no vean; paso por alto que aportó una exposición tan ridícula inventada por él. Pero vendré a otro sentido suyo, donde dijo muchas cosas; mas yo, no atendiendo a la singularidad, que es madre de los errores, sino siguiendo la sentencia de los santos padres y de la santa madre Iglesia recibí la palabra «Apacienta», porque allí lo constituyó pastor y le dio el primado de la Iglesia universal. Así lo entendieron Gregorio, así Crisóstomo, así Ambrosio sobre el último de Lc, ayer por mí no alegado. A los cuales pasando bellamente 'de largo' nada dijo el reverendo 'padre, sino que aportó la verdadera sentencia de Agustín, de cómo debe ser el prelado o pastor, la cual también yo admito, y creo que cualquier fiel ruega que también el sumo pontífice, 'cualquier príncipe', cualquier prelado buscase el honor de Dios y la salud de los súbditos, pues este es el verdadero apacentamiento, por el cual, así como aquí tienen el principado en la tierra, también son elevados para que reinen en los cielos. Por lo cual Agustín no dice cosas adversas ni pugnan­tes, sino tales que por ellas debe entender el prelado cuál sea su institución.

Y no vale lo que el reverendo padre adujo de que le fueron encomendadas las ovejas bajo condición, no habiendo allí condición alguna, salvo que su paternidad tenga otro evangelio que yo. Mas significando cómo debe ser el pastor, Cristo busca en Pedro la dilección y el amor de Dios, y así lo hizo pastor de toda la Iglesia, como canta la santa madre Iglesia, y el reverendo \marginpar{\small[308]} padre ora: 'Tú' cabeza de la Iglesia, pastor de las ovejas. Por lo cual tenemos por derecho divino en este lugar que el prelado debe amar a Dios, tenemos a Pedro constituido pastor de las ovejas, como lo entendieron san Gregorio, san Ambrosio y Crisóstomo, aun el griego.

Nada importa la tenue racioncilla aducida en contrario por el reverendo padre, que Pedro no apacentó a todos; lo que se hace en nuestro nombre se dice que lo hicimos nosotros, como atestigua Clemente que Pedro ordenó obispos, presbíteros y diáconos que sembrasen el Evangelio por todo el orbe de la tierra, como se lee en el can. In illis, dist. 80 o 81. De donde el enviado por el principal, como fue san Pablo, no está obligado a ser enviado por el vicario del mismo; sin embargo Pablo subió con Tito, a Gal 2, y confirió con Pedro el Evangelio, quienes le dieron la diestra en señal de compañía.

Además, en cuanto a que forja contra mí una glosa jamás por mí pensada, como si yo hubiese dicho: Edificar sobre la piedra es estar sometido al Romano pontífice, ¡oh sueños, oh portentos de palabras!; jamás fue esto de mi mente. Pues Cristo edifica la Iglesia, no los que están sometidos, ya que los súbditos deben estar sometidos a la Iglesia edificada.

Sobre Gregorio de Rímini, cuya sentencia me desplace y a él le place, nada ya al propósito. No quiero, pues, descender a estos desvíos. Sino

como hemos recibido, defendamos que el artículo del concilio de Constanza es verdadero, y que quienquiera que lo contradiga yerra y contradice a la verdad y al honor del concilio.

\paragraph*{El D. D. Martín.}

Antes de venir a las autoridades por mí opuestas, recorreré brevemente lo que el egregio D. D. ha mezclado sobre mi respuesta de ayer y de hoy, donde entre otras cosas, como desesperado tras consumir todas las armas, coloca en un solo artículo del concilio de Constanza el postrer presidio de su confianza, y odiosísimamente lo repite e inculca, como si arriba no se hubiese respondido a él bastante sobradamente. Diré, sin embargo, una vez más.

Cuando el artículo sobre las obras de los hombres se cuenta entre los condenados y es afirmado por san Pablo y Agustín, defendido después por Gregorio de Rímini y por todas las universidades hasta hoy, no me moverán las odiosísimas inculcaciones de este artículo, hasta que el egregio D. D. pruebe que el concilio no puede errar, que no erró o aun que no yerra, pues el concilio no puede hacer derecho divino de lo que por su naturaleza no es derecho divino; por eso tampoco es herético sino lo que es contra el derecho divino. Esto por este artículo.

En cuanto a que zahiere que Crisóstomo fue legítimamente depuesto, no me entendió, pues claramente dije hoy que la potestad de ordenar y deponer al obispo fue legítima en Epifanio de Chipre y en Teófilo, aunque nada disputé sobre si entonces usaron bien o mal esta potestad legítima. Me basta que, sin \marginpar{\small[309]} la autoridad del Romano pontífice, tantos santos y celebérrimos varones de Grecia, catolicísimos, no cismáticos, tuvieron potestad de ordenar y deponer obispos.

También pondera egregiamente que yo habría acusado de ambición a los santos antiguos, y no atiende a que dije que también en los apóstoles estuvo esta fragilidad en presencia de Cristo.

En cuanto a que ahora engrandece tanto a Platina contra san Gregorio, por mi autoridad consta bastante claro a cualquiera aun medianamente ingenioso qué concluye. Pues yo de Platina tomo las historias y nada más. Y que el mismo Platina aduce ciertas razones, lo rechazo por la autoridad contraria de Gregorio.

Ya paso por alto aquello de que Sixto y Víctor escribieron llamándose pontífices universales y que a León fue aclamado en el concilio Calcedonense, sobre lo cual

hoy se ha dicho bastante. Hay que observar aquello de que, según dice el egregio D. D., los pontífices siguientes, siguiendo en Cristo la forma de la humildad, se abstuvieron de este nombre; lo cual no puedo entender de otro modo sino que por causa de humildad detrajeron al derecho divino, y se hicieron mendaces y desobedientes, si es derecho divino, esto es precepto, que ellos fuesen pontífices universales.

Repitiendo la carta responsiva de Inocencio a Agustín dice que fue antigua costumbre consultar al Romano pontífice en lo que es de la fe. Óptimamente dice, y a mí me place, sobre todo en la Iglesia latina; pero ¿qué tiene que ver con la fe y el primado? Esto es, ¿a qué propósito?

También se quejó de que yo exigí de él la probación de la solución sobre los veinte años; doblemente falta el egregio D. D.: primero, porque él mismo fue oponente durante tres días y medio y solo a esta hora empezó a ser respondiente. Después, no exigí que probase la solución, sino el primado trasladado con Pedro.

Excusándose también sobre Cefas —que Pedro se interpreta en Juan—, aún persiste en que Pedro también significa cabeza de la casa, como si, a saber, Juan, para interpretar el Evangelio a los griegos, hubiese hablado no en griego sino en caldeo, para no ser entendido acaso.

La palabra «apacienta» la toma de modo que quiere que signifique la constitución del pastor universal, y para ello adujo a Gregorio, a Ambrosio y a Crisóstomo, ayer no alegados, sobre los cuales nada diría yo. Diré, pues, ahora, si es que fui olvidadizo, y a la vez opongo la regla de Agustín: Que los dichos de todos los escritores han de juzgarse por la divina Escritura, cuya autoridad es mayor que la capacidad de todo el género humano; no porque condene la sentencia de los ilustrísimos padres, sino porque imito a los que más de cerca se acercan a la Escritura, y ante todo, si la Escritura fuese clara, abrazo la Escritura misma. Siendo, pues, claro que apacentar, por fuerza del vocablo y por uso de la Escritura, no significa sino enseñar y cuidar \marginpar{\small[310]} las ovejas, como Pedro lo recibe al final: Apacentad la grey que hay en vosotros [1 Pe 5 v. 2]. No conviene equivocar el vocablo y dejar sin necesidad el sentido simple.

También disuelve lo que dije que hay condición en la palabra de Cristo, Si me amas, cuando yo juzgo que no solo hay condición, sino también pacto

entre Cristo y Pedro, cuando aquél dice: ¿Me amas? y éste responde: Te amo.

Omito la dialéctica, en que ejercitadísimo es el señor D. D., y la conjunción condicional «si». Querría oír qué diría a la decretal de elect. c. Significasti, donde el pontífice Pascual dice manifiestamente que aquí hay condición, si al menos tanto vale la autoridad del Romano pontífice como hasta aquí ha contendido; dice en efecto. [Va a citar del libro esta decretal.]

A mi razón, que así la confutó —que Pedro no apacentó a todos, porque, dice, lo que se hace en nuestro nombre se dice que lo hicimos nosotros—. Aunque no creo que con esto haya querido el señor D. satisfacer a mi razón, pues no puede decir que Pablo predicó en nombre de Pedro, quien en cambio a Rom 1 [v. 1ss.] se dice apóstol en nombre de Cristo; lo cual colijo de que dice que Pedro, con Clemente por autor, ordenó obispos, presbíteros y diáconos que sembrasen el Evangelio por todo el orbe, entre los cuales no cuenta a Pablo.

Pero esto parece hacer al caso, que el enviado por el principal no está obligado a ser enviado por el vicario del mismo; esto admito, si primero fuese probado que Pedro es vicario. Ahora, en cambio, no menos es Pablo vicario de Cristo que Pedro, como ya se ha dicho, porque es apóstol en nombre de Cristo, y así tantos vicarios cuantos apóstoles.

A propósito.

A la autoridad de Pablo 1 Cor 3 [v. 5] dice primero que Cefas no es del texto desde el principio, sino que se recuerda al final, donde la glosa confiesa que Pedro es mayor que todos; cosa levísima es esto, pase, pues al principio del primer capítulo [v. 12] ciertamente se pone Cefas.

Segundo, nada dice sino que ni palabra ni ápice hay en el texto sobre la negación del primado, y así quiere que esta autoridad quede disuelta. Añadió, sin embargo, la autoridad de Jerónimo, Agustín, Ambrosio, Crisóstomo, Bernardo y León, que dicen que Pedro es la piedra. Contra lo cual 'digo y' replico que este segundo dicho nada es al propósito y arriba está bastante disputado si Pedro es la piedra. Igualmente tampoco disuelve aquello de que dice que no hay palabra ni ápice sobre la negación del primado; pues si aun Jerónimo establece la cabeza para que se quite la ocasión del cisma,

bastante robustamente 'lucha en contra' el texto abierto, que quita la cabeza para que se quite la ocasión del cisma. Dice en efecto: ¿Qué es Pablo? ¿Qué Apolo?, queriendo que solo Cristo sea la cabeza, con lo cual verdaderamente se quita todo cisma. Y así no se satisface con esta respuesta.

\marginpar{\small[311]} A la segunda autoridad, Gal 2, dice que ciertamente todos los apóstoles fueron iguales, como atestiguan Jerónimo y el can. In novo. Pero disolver tan gran autoridad con estas tres palabras —que otra cosa es la potestad de régimen y de administración— me es bastante admirable. Pero tampoco entiendo qué quiere, cómo distingue el régimen y la administración.

Y en cuanto a lo que dijo, que Pablo 'con Tito' subió a Pedro para conferir el Evangelio con ellos, claro es por el texto precedente y siguiente, como el mismo Jerónimo expone, que Pablo no hizo esto para consultar a Pedro sobre el Evangelio, que ya por más de 14 años había predicado y del que se había jactado de no haberlo recibido de hombre, sino para cerrar la boca a los pseudoapóstoles, para confirmar a los Gálatas que eran lo mismo lo que él y los demás apóstoles predicasen.

\paragraph*{El D. D. Eck.}

En cuanto a que el reverendo padre mezcla el artículo de Juan Hus sobre la obediencia eclesiástica con otro artículo sobre las obras de los actos humanos, por autoridad de Gregorio de Rímini, que sería defendido por todas las universidades, al cual hasta ahora en ese lugar en ninguna universidad, aunque en muchísimas estuve, oí defender; pide, sin embargo, de mí que pruebe que el concilio no puede errar; no sé qué quiere esta petición, si tácitamente quiere tener por sospechoso al laudable y glorioso concilio de Constanza. Esto os digo, reverendo padre: si creéis que el concilio legítimamente congregado erró o yerra, sois para mí como étnico y publicano. Qué sea hereje, por ahora no lo discutamos.

Sobre Crisóstomo, no quiere haber justificado su expulsión, sino haber alabado la ordenación de obispos instituida por el concilio Niceno. Me place por reverencia a san Crisóstomo, y siempre los Romanos pontífices en diversas

permitieron a las naciones diversos modos de elegir a los obispos, consta además con suficiente claridad, a partir de Gregorio Nacianceno y de Crisóstomo restituido, que intervino la autoridad del pontífice romano, como atestiguan san Atanasio de Alejandría, san Pablo y otros obispos de la Iglesia oriental, que recurrieron al pontífice romano para su restitución. Que también en los apóstoles hubo algo de humana fragilidad, lo sabemos. Pero mientras desempeñaban públicamente su oficio, sin indicio ni causa no ha de detraerse a tan santos varones.

Acerca de Platina dice que acepta la autoridad de su historia, y yo tampoco aduje razones tomadas de Platina, sino la historia.

A la carta de Inocencio responde que nada atañe al propósito, porque prueba acerca de la Iglesia latina, no acerca de otra. Pero como esta es una glosa luterana sin apoyo en autoridad ni persuasión alguna, la desprecio con la misma facilidad con que se aduce.

A lo de los veinte años respondió que yo había sido oponente ya por tres días; y ciertamente debería haberlo sido, si el reverendo padre con sus cautelas no hubiese convertido alguna vez de opo\marginpar{\small[312]}nente en respondiente. Pues este argumento no lo puse yo, sino que lo puso él mismo, y ciertamente aquella cátedra Pedro no se la llevó consigo a Roma, sino que trasladó el primado.

Sobre la palabra «apacienta» quiere dejar de lado a Gregorio, a Ambrosio, a Crisóstomo, porque una palabra simple no debe tomarse en sentido equívoco. Digo: ¿quién ignora que la palabra pastor es nombre de oficio, de honor y también de carga? Por eso recibió a la vez el oficio de la prelacía y juntamente la carga de apacentar.

La condición aducida por el reverendo padre no la acepto, porque tampoco Cristo, como él mismo siempre se adhiere a la letra, unió condición alguna al apacentar. Pues así se confirmaría el error husita de que quien no está en caridad no sería obispo, prelado, pastor.

Sobre el capítulo Significasti, de elect., sé que Gerson y otros doctores discuten en la materia si el concilio está por encima del papa. Pero como no se han aducido las palabras de la decretal, nada puedo decir, sino que conjeturo

«sin querer ligarme a esto», que aceptó la condición pascual por cualidad.

Sobre Pablo, que no reconoció el primado de Pedro, sino que fue apóstol de Cristo: todos somos de Cristo, creo; sin embargo, el primado de régimen a Pedro en ninguna parte se lo quitó, sino que incluso al morir, como se escribe en la carta sobre el tránsito de los apóstoles, se despidió de Pedro como de cabeza de la Iglesia; carta de la que no tengo certeza si es del bienaventurado Dionisio o no.

Dijo que Pablo fue vicario de Cristo, pero no general, pues también en el Apocalipsis [21, v. 14] se ponen doce fundamentos de la Iglesia, y sin embargo otro fundamento, a saber, el principal, nadie puede ponerlo sino el que ha sido puesto, es decir, Cristo.

A lo segundo dijo que no viene al propósito si Pedro es la piedra, y ese es todo nuestro asunto. Sin embargo, silenció la solución dada por mí anteayer, a la que me remití, aunque está expreso en la letra que Pablo reprueba las contiendas. Pero no sé con qué ingenio el reverendo padre juzga que Pablo negó por ello el primado, para que se quitara el cisma. Quizá no estuvo en una provincia que careció de príncipe, pues habría experimentado lo contrario: cuántas disensiones, contiendas y discordias surgen donde no hay príncipe ni cabeza que juzgue y conozca entre las causas.

Y le parece admirable que con tan pocas palabras, esto es, apostolado y régimen, quiera yo resolver tan grande autoridad. Los principios son mínimos en cantidad, máximos en virtud. Si el reverendo padre hubiese atendido a aquellas dos menudas palabras, apostolado y régimen, de la cabeza principal y de la secundaria, hace tiempo habría reconciliado las sentencias de la Sagrada Escritura y de los santos padres que le parecen pugnar entre sí; descuidadas estas, es necesario seguir dogmas exóticos y ajenos.

Mañana a las siete de la mañana se continuará la disputa por edicto público.

\marginpar{\small[313]} El viernes ocho del mes de julio a las siete de la mañana.

\paragraph*{El D. D. Martín.}

A las autoridades aducidas por mí, de Pablo a los Corintios y a los Gálatas, el egregio D. D. nada responde a su modo, sino que repite lo del concilio de Constanza y los husitas, con lo cual no se me resuelven las autoridades. Pues, como dije, el concilio de Nicea estableció otras cosas mucho más sagradas y célebres. Por eso, cuanto él se apoya en el suyo, tanto yo en el mío, y digo, como dije ayer: por mucho concilio que se alegue, aún no se tiene el derecho divino, del cual solo se trata; luego el concilio nada viene al propósito.

En segundo lugar, para eludir las autoridades más que para resolverlas, trajo la distinción entre apostolado y administración, la cual, puesto que ha sido hallada por arbitrio humano y no se encuentra ni en las sagradas letras ni en los padres eclesiásticos que las interpretan, no me doy por satisfecho con que se me haya devuelto, sobre todo cuando apostolado no es, como él piensa con los suyos, nombre de dignidad, sino de oficio, como es claro por el apóstol a los Romanos [1, v. 5] y en todas partes, donde dice que recibió la gracia y el apostolado para la obediencia de la fe. Pues apóstol es mensajero de la palabra, y apostolado o ἀποστολή en griego es oficio de la palabra, y así en las sagradas letras apostolado y administración serán casi lo mismo; a no ser que entienda por administración las contiendas forenses sobre bienes temporales eclesiásticos, la cual no es propiamente administración apostólica. Confieso, ciertamente, que si se permite esta libertad de fingir, fingiré también yo que san Juan es canciller y que otro apóstol es otra cosa. Pero todo esto no es de derecho divino. Se desvanece, pues, aquel principio que jactó de mínimo en cantidad. Y bastante me admiro de que el egregio señor rehúya así las sagradas letras y huya de su rostro, desviándose a estas galerías inventadas por ingenio humano.

Por tanto, sigue aún en pie la autoridad a los Corintios, según la cual, para quitar el cisma, Pablo quita de en medio la persona de Pedro, de Pablo, de Apolo, etc., y deja solo a Cristo, quedándose para sí únicamente el ministerio de la palabra, que es verdaderamente el apostolado y la administración de la Iglesia. Y no vale lo que con aquella fría semejanza quiso disolver esta robustísima autoridad,

a saber, que una provincia sin cabeza queda expuesta a cismas y a muchos peligros, pues no faltará Iglesia sin cabeza cuando se den muchos obispos sin papa.

Como vemos que sucede en las federaciones, donde tanto más hermosamente se administran las repúblicas cuanto menos tienen una cabeza común, pues es evidente que el reino de Francia, el de España, el de Hungría y otros infinitos en el mundo existen sin cabeza común y sin embargo no chocan entre sí, contentándose cada reino con su rey. Y si en estas cosas temporales y fluctuantes, no teniendo los hombres, tan diversos en ingenios y costumbres, cabeza común ni en el cielo ni en la tierra, se mantienen la paz y la concordia de las tierras, cuánto \marginpar{\small[314]} más podrá suceder esto en las cosas espirituales, constantísimas, esto es, en hombres dotados de una misma fe, de igual caridad y de las mismas costumbres, y que además tienen en el cielo una cabeza común, como dice el apóstol: Una fe, un Señor, a los Efesios [4, v. 5]. Por lo cual, aunque hubiese diez papas, aunque hubiese mil papas, no por ello se producirían cismas.

Por eso pido al D. D. que quiera resolver mejor aún la autoridad de Pablo, y asimismo aquello a los Gálatas [2, v. 6]: Dios no hace acepción de personas, donde Pablo rechaza todo lo que en Pedro y en los apóstoles puede haber de personal o cualificado, como dicen. Aún no ha sido tocada, menos aún resuelta.

A lo cual, para terminar la disputa, añado también aquello de que Pablo, en 1 Corintios 12 [v. 28], al describir el cuerpo eclesiástico, enumera en primer lugar a los apóstoles, en segundo a los profetas, en tercero a los doctores y casi en último lugar a los gobiernos, que el D. D. llama administraciones; luego, si el primado de Pedro fuese cosa tan grande y tan necesaria por derecho divino, Pablo habría descrito insuficientemente la Iglesia, al omitir lo principalísimo que hay en ella.

Además, sigue aún en pie esto: que ni Pedro ni todos los apóstoles juntos pudieron ordenar a Matías, Hch 1, ni darle la administración. Por eso el pontífice romano se arroga más administración para conferir por todo el mundo de la que no se encuentra en su predecesor el apóstol.

Igualmente pienso acerca de la ordenación de Pablo y Bernabé, Hch 13, a los cuales el Espíritu Santo, y no Pedro o la Iglesia, segregó. Por eso se gloría no sin razón de no haber recibido nada de los grandes apóstoles; más aún, escribiendo a los Corintios dice que en nada hizo menos que los demás apóstoles [2 Cor 12, v. 11].

Pero también aquella autoridad estrecha aún más a Pedro, en la que escribe a los Gálatas: El que obró en Pedro para el apostolado de la circuncisión, obró

también en mí entre los gentiles, y más abajo: Me dieron a mí y a Bernabé las diestras en señal de comunión, para que nosotros fuésemos a los gentiles y ellos a la circuncisión, con tal de que nos acordásemos de los pobres [2, vv. 8 ss.]. Si quisiera apoyarme en esta autoridad, ni el D. D. ni hombre alguno podría obligarme a que el apostolado de Pedro sea más amplio que a los judíos por derecho divino, y por ello todas las autoridades en que se apoya —Apacienta mis ovejas, Tú eres Pedro, etc.— han de referirse a los términos de su apostolado, como Cristo claramente en Jn 10 [v. 16] separa las ovejas de los gentiles de las ovejas de Israel, diciendo: Tengo aún otras ovejas que no son de este redil.

Y al fin, como también ayer tocó el D. D., en Ap 21 [v. 14] se describe la nueva Jerusalén con doce fundamentos, además de aquel fundamento capital que es Cristo, y sin embargo allí no se pone ninguna diferencia entre los doce fundamentos, lo cual habría que hacer del todo si hubiera de pensarse algo tan grande acerca del primado de Pedro. A lo mismo valen muchísimas figuras en las Escrituras, como los doce bueyes del mar de bronce, 3 Re 7, asimismo los doce leoncillos del trono de Salomón, las doce piedras del Jordán y otras semejantes, en todas las cuales se describe una total igualdad de todos los \marginpar{\small[315]} apóstoles, y ninguna desigualdad puede probarse en contrario por derecho divino.

Sea esto dicho en favor de esta conclusión, que hasta hoy creo verísima, más aún, sé con certeza que no podrá ser nunca demolida. Por eso ruego al egregio D. D., puesto que aún no ha opuesto ninguna autoridad de derecho divino en favor de su sentencia, que al menos disuelva algunas clara y sin rodeos; de otro modo no lo tendré por teólogo que siente rectamente acerca de las Escrituras, y lo dejo, o más bien lo remito, al juicio no de la envidia ni del vulgo, sino de los jueces que han de designarse y de cualquier lector prudente y de todo oyente sincero.

\paragraph*{El D. D. Eck.}

No sé qué quiere el reverendo padre, que al inicio de su discurso dice que nada he respondido, y sin embargo en el curso del mismo se esfuerza por quebrantar las soluciones dadas por mí ayer, invictas y robustísimas; si esto es coherente, él mismo lo verá.

Después, en cuanto a que no quiere verse constreñido por la autoridad del concilio de Constanza, él verá, vean los jueces cuán justamente obre. El sínodo niceno

no lo rechazo en lo que es de fe, pero lo que mira a las costumbres, como las ordenaciones de los obispos, puede cambiarse según la calidad de los tiempos y lugares.

Me pide que le aporte el derecho divino, lo cual, aunque ya lo he hecho muchas veces, sin embargo mi conclusión no lo contiene, sino solo esto: que el reverendo padre dijo en su resolutorio que la Iglesia romana antes de los tiempos de Silvestre no fue superior a las demás, lo cual niego en mi conclusión; con todo confieso que el primado del sumo pontífice es de derecho divino según lo alegado por mí.

Después, para reprobar mi solución dice que mi distinción ha sido inventada por arbitrio humano, como si en teología no fuese lícito distinguir, igual que Arrio quiso burlarse de Atanasio, pues, demasiado aferrado a la letra, dijo que el Hijo es criatura por aquello: El Padre es mayor que yo [Jn 14, v. 28]. Al disolverlo Atanasio distinguiendo: El Padre es mayor que yo según la humanidad, pero no según la divinidad; como la Sagrada Escritura no había distinguido, Arrio no la aceptó y así permaneció pertinazmente en el error. Así en el presente caso, cuando los que esclarecen la Sagrada Escritura tendrán vida eterna, hay que dar distinciones cotejando los diversos pasajes, para que no se admita contrariedad en las sagradas letras, no sea que, por aferrarse demasiado a la letra, uno sea muerto por ella, sino más bien vivificado por el espíritu. Y como Pedro es antepuesto a los demás muchas veces en las sagradas Escrituras y muchas veces se le encuentra igual a ellos, para evitar la contradicción es necesario poner una distinción conforme a la Sagrada Escritura.

Así también lo insinúan Jerónimo y Cipriano por mí alegados, quienes confiesan la igualdad del apostolado y sin embargo tributan a Pedro el primado, lo cual no puede hacerse si no se refieren a cosas diversas. \marginpar{\small[316]} Me place decir que el apostolado es oficio, pero de modo que también diga honor. Pues estos nombres son de honor y de carga, como dije, según indica el apóstol a los Romanos, cap. 10 [v. 15]: ¿Cómo evangelizarán si no son enviados?

Sobre la administración, de nuevo doy mi asentimiento al reverendo padre, a saber, que aquella en lo temporal es muy inferior a aquella en lo espiritual. Roguemos al Dios óptimo que nuestros obispos reconozcan esto y administren lo espiritual no por vicarios y lo temporal por sí mismos. Llamé, pues, administración a la del régimen en lo espiritual de modo potestativo, y no obsta lo que dice el reverendo padre, que así él podría fingir que Juan es canciller, Marcos portador del sello, etc., porque, como dice Agustín, los filósofos usan palabras libres, pero a nosotros solo nos es lícito hablar según una regla cierta; por eso conviene aportar cosas que reciban fundamento de la Sagrada Escritura,

como el primer miembro, el apostolado, en la igualdad de los apóstoles, y el otro miembro, el régimen, en el primado de Pedro. Por eso confieso que la autoridad de Pablo es robustísima, pero que en ni una sola letra me es contraria.

Despreció el reverendo padre el símil aducido por mí sobre los cismas en la provincia carente de gobernante, porque ya sucede así en los reinos fieles de España, Francia, Inglaterra y Hungría, que existen sin cabeza, y sin embargo se mantienen entre ellos la paz y la concordia. Me admiro de cómo el reverendo padre ha podido olvidar tantas matanzas, guerras y batallas, que no puede ignorar. No quiero decir nada de tan diuturna desavenencia y perpetuo odio entre franceses e ingleses; ni del odio inveterado entre franceses y españoles, que tantas veces por el reino de Nápoles libraron guerras con máxima efusión de sangre cristiana, y ojalá al fin la paz fuese estable y permanente. Reconozco con el apóstol una sola fe, un solo Señor Cristo, pero venero al pontífice romano como vicario de Cristo. Queda, pues, la primera solución invicta y válida, al menos a mi exiguo juicio, del cual juzguen los más doctos.

Además, el padre digno de respeto dijo que las palabras de Pablo no han sido tocadas por mí: Dios no hace acepción de personas, y por tanto menos aún resueltas. Quizá oyó mal, pues tanto por Ambrosio como por Jerónimo aduje el sentido de esas mismas palabras: Juan y Pedro habían sido idiotas e iletrados en la ley; Pablo, en cambio, había sido doctísimo, aventajando a muchos de sus coetáneos, como de sí atestigua [Gá 1, v. 14]. No quiso, pues, Pablo estimar lo que habían sido Pedro y Juan, porque Dios no hace acepción de la persona del hombre, pues así eligió al pescador como al docto escriba. Pero con ello no se quita el primado, ni se imputa a Dios acepción de personas, aunque haya preelegido a uno sobre otro, como Pedro confesaba de Cornelio [Hch 10, v. 34]: En verdad he comprendido que no hay acepción de personas ante Dios.

Vengamos a lo aducido por el reverendo padre. \marginpar{\small[317]} Cita al apóstol, 1 Cor 12 [v. 28], que al describir el cuerpo místico refiere apóstoles, profetas, doctores y en último lugar la administración. Digo que el apóstol describió egregiamente, pero que allí ni negó ni afirmó el primado; y es inválido el argumento: aquí Pablo no recuerda tal primado, luego el primado no existe en la Iglesia, tanto porque hasta los niños saben que el argumento por autoridad negativa no vale, como porque también lo no expresado en un lugar de la Sagrada Escritura debe tomarse, según la doctrina de Agustín en el libro sobre la doctrina cristiana, de otro lugar, como hemos hecho nosotros.

De igual modo respondemos a aquello de Ap 21. Confieso que son doce los fundamentos, y que estén figurados también por los doce bueyes, por los doce leoncillos, por las doce piedras del Jordán; pero en ellos se afirma el número duodenario de los apóstoles, y no por ello se niega el primado. Yo, en cambio, aduje aquel lugar porque es preciso mostrar concordes los pasajes de la Sagrada Escritura, y habiendo puesto el apóstol un solo fundamento, Cristo, y no pudiendo nadie poner otro, no debe negarse por ello que Pedro sea fundamento de la Iglesia, porque también el mismo san Juan puso doce fundamentos; lo cual, si no aceptas la distinción entre fundamento de los fundamentos, esto es, principal, y fundamento sustituido o secundario, no concordarás a Pablo con Juan.

En segundo lugar adujo que Pedro no pudo ordenar a Matías, y que así el Espíritu Santo segregó a Pablo y Bernabé. Respondo: como los apóstoles fueron obispos, no recuerdo haber encontrado en la Sagrada Escritura que fuesen ordenados obispos por Cristo, aunque en la última cena los ordenó primeramente sacerdotes. Por eso también puedo referir aquello a Pedro en cuanto a la ordenación episcopal; sin embargo, como admito que el oficio del apostolado solo es tributado por Dios, no es extraño que pidieran el auxilio divino, pues echaron suertes no al modo vulgar, sino que, como atestigua el sagrado Dionisio, aquella suerte fue una visible aparición del Espíritu Santo. Por eso aquel lugar habla en favor del primado de Pedro, porque ejecutó el primado que se le había concedido: En aquellos días, levantándose Pedro en medio de los hermanos, etc. Como también en lo de Safira y Ananías, como en la excusa de los apóstoles de que no estaban llenos de mosto, como en la respuesta y defensa ante el concilio de los judíos [Hch 1-5] y en otras cosas contenidas en los Hechos de los Apóstoles, donde ejerció el régimen de la superioridad. Por eso creo que Matías fue hecho apóstol por el Señor y ordenado obispo por san Pedro. Así también hace tiempo concedí de Pablo que recibió de Dios el apostolado, y confieso que trabajó más abundantemente que los otros.

Pero vengamos a lo tercero, que entre lo aducido por el reverendo padre parece hacer más al propósito: El que cooperó con Pedro, a los Gálatas 2 [v. 8]. Y digo que Pablo no promulgó allí un derecho divino, sino que narró un hecho, pues también Pedro en los Hechos de los Apóstoles fue mandado instruir en la fe a Cornelio el gentil, y a los gentiles en Antioquía, donde presidió, e igualmente en Roma, donde al fin fue coronado con el martirio por Nerón el gentil, presidió y enseñó juntamente a los judíos convertidos y a los que habían \marginpar{\small[318]} de convertirse a la fe.

Por eso consta claramente que Pablo narró un hecho. Paso por alto que san Pablo no recuerda el oficio del régimen o del primado de la presidencia, sino la ejecución del apostolado, según la cual los apóstoles se repartían por suerte diversas provincias, en lo cual también Pedro fue igual a los demás apóstoles según Anacleto, can. In novo, según Jerónimo contra Joviniano y sobre la carta a los Gálatas, según Cipriano en la carta 9 a Cornelio papa, lib. 4. Por esta razón la Sagrada Escritura ha de entenderse con efecto, de modo que advirtamos cuándo habla del oficio del apostolado y cuándo del régimen de la potestad o presidencia.

No importa que el reverendo padre niegue aún que yo haya aducido derecho divino en favor del primado y que no haya disuelto las autoridades aducidas por él, porque lo mismo diría yo de él. Sin embargo, ninguno de los santos padres entendió las autoridades aducidas por el reverendo padre para la negación del primado de Pedro sobre los demás, como él las adujo. En cambio, de múltiples maneras los santos padres entendieron que las autoridades aducidas por mí hablan del primado de Pedro. Por eso, en torno a esta conclusión, si alguna vez hablé algo con más acritud y dureza contra el reverendo padre, no quisiera que pensarais que fue dicho contra su persona, sino que la causa misma que se ventila me parecía tan preñada y necesaria que alguna vez hubieron de usarse palabras más acres. En las demás conclusiones experimentaréis siempre la modestia eckiana.

Por eso también yo quiero terminar por ahora esta conclusión, porque reputo que san Pedro tuvo de Cristo el primado de toda la Iglesia por la promesa a él hecha en Mt 16, como entienden Jerónimo, Cipriano, Agustín y otros; por la nominación, por la que es nombrado primero en Mt 10, según la glosa allí mismo; por el pago del tributo, Mt 17, donde solo él fue equiparado a Cristo, según Agustín y Ambrosio; por que Cristo rogó por su fe indefectible y le mandó confirmar a los hermanos, Lc 22, según Crisóstomo y san León; porque fue constituido pastor de las ovejas de Cristo, Jn último [cap. 21], según Crisóstomo y Gregorio; porque Cristo dice a Pedro: Sígueme, no solo en el género del martirio, sino también en el orden del magisterio, según Teófilo; porque solo Pedro recibió el mandato de venir sobre el mar hacia Cristo [Mt 14], para que así el mundo, designado por el mar, le estuviese todo sometido según Bernardo en su consideración a Eugenio, y por otras cosas aducidas arriba, reputo verdadera la conclusión comprobada por tantos santos padres y últimamente por el concilio de Constanza. No quiero, sin embargo, apoyarme en mi propia prudencia, dispuesto

a cautivar mi entendimiento en estas y en cualesquiera otras cosas al juicio de los jueces que han de designarse, de la sede apostólica y de otros varones prudentes y buenos.

Solo a Dios la gloria. Eck.

A las dos de la tarde se continuará la disputa.

\marginpar{\small[319]} A las dos de la tarde del mismo día.

\paragraph*{El D. D. Martín.}

Nos vemos obligados a añadir un remate, porque el egregio D. D., al disolver mis objeciones, opuso a la vez muchísimas cosas y de nuevo me obligó a responder. Y aunque hubiera querido pasar esto por alto, sin embargo, para declarar plenamente mi sentencia a los jueces que han de designarse, respondo en pocas palabras.

Primero objetaba que Pedro presidía la elección de Matías, exhortando a los apóstoles, etc. Creo, sin embargo, que ni al mismo D. D. le parece que por esto se pruebe el primado de Pedro, si Pedro predica, pues en lo anterior él mismo me respondió que una cosa es el apostolado y otra la administración del régimen. No dijo, pues, rectamente que Pedro al predicar así ejerció la administración del régimen; de otro modo también Esteban será el primero y papa, porque también él predicó e hizo muchos signos en el pueblo.

Y en cuanto a que cree que Matías fue hecho apóstol y ordenado obispo por Pedro, me doy por contento con que crea lo que quiera; yo no lo creo si no lo prueba. Asimismo, que Pedro ejerció el primado en el castigo de Ananías y Safira puede decirse, pero no se prueba, pues también resucitó a muertos, lo cual también hicieron otros apóstoles. Asimismo, que excusó a los apóstoles para que no pareciesen ebrios tampoco es indicio de primado, pues, como dije, el egregio D. D. distinguió siempre el primado de todas las obras y administraciones comunes a los demás. Asimismo, que Pedro caminó sobre el mar y que Bernardo interpreta por el mar el mundo, lo concedo, pero nada atañe al primado, pues también nosotros debemos hollar el mundo.

Aquello de la indefectibilidad de la fe prometida por Cristo a ruego suyo lo paso de buen grado, pues concedo que la fe de Pedro nunca cayó, aunque él mismo cayó.

por la fe. Pues entonces el ladrón creyó con la fe de Pedro, cuando Pedro negó, como dice Agustín. Muy otra cosa es la fe que el primado.

Asimismo, en cuanto a aquello de que Pedro pagó el estáter por Cristo y por sí, en lo cual se dice que fue igualado a Cristo, lo concedo, pero nada aporta al primado, antes bien va contra el primado, pues enseguida sigue, como también interpreta Jerónimo, que los demás apóstoles, ofendidos por sospechar que Pedro habría de ser mayor, comenzaron a discutir quién sería el mayor. Por eso Cristo, llamando junto a sí a un niño, reprimió su ambición, y así lo que sigue del texto combate el primado con más fuerza que lo que precede lo favorece.

Y en cuanto a aquello del último capítulo de Juan [c. 21], que lo dicho sólo a Pedro: Sígueme, según Teófilo haya de entenderse no sólo del género de martirio, sino también del orden del magisterio, rechazo a Teófilo, tanto porque Agustín parece sentir mejor al decir así: Sígueme padeciendo males temporales. Y aunque Agustín no lo dijera, el propio texto del Evangelio lo probaría. Pues corrió la voz entre los discípulos de que aquel discípulo no moriría, al oír decir a Cristo: Si quiero que él permanezca, ¿a ti qué? Tú sígueme. Por eso es manifiestísimo que aquella imitación ha de entenderse del género de martirio y de la pasión.

\marginpar{\small[320]} Por lo cual ruego que el D. D. omita traer cosas nuevas, y sobre todo las que no hacen al caso, cuando también yo podría aducir aquello de Hch 8 [v. 14], donde los apóstoles enviaron a Pedro y a Juan como a un inferior, y Hch 15 [v. 13 ff.], donde Santiago confirmó y modificó el discurso de Pedro. Pero no he querido aducir estas cosas ni otras semejantes, pues hace ya tiempo que concedí a Pedro sólo el primado de honor y únicamente negué el primado de potestad sobre los demás apóstoles. Pues ésta es aquella acepción de persona, más aún, aquel primado de honor, que Dios no tiene en cuenta, como arriba aduje, y no aquella persona del pescador o de la pobreza, como expuso Ambrosio según refiere el doctor, pues Jerónimo interpreta mejor a Pablo por la magnífica cualidad por la que los pseudoapóstoles quisieron subvertir a los gálatas como apartándolos de la autoridad de Pedro. Estoy pues muy contento de que el D. D. haya confirmado su proposición con autoridades de Jerónimo, Ambrosio, Bernardo y León, sobre todo a partir de esta autoridad: Tú eres Pedro, con tal de que mi sentencia, como se ha dicho en lo precedente, se tenga por más confirmada con autoridades iguales y mayores y que hablan más a la intención, remitiendo también esto al juicio, como arriba.

\paragraph*{El D. D. Eck.}

El reverendo padre vuelve a abrir hoy lo ya concluido y rebate como principales las cosas que yo había añadido como accesorias y adornos en la coronación de la conclusión, y hasta tal punto procura la metamorfosis, que ignora si se ha transformado en oponente o en respondiente. A la cosa, brevemente.

Nunca soñé aquella consecuencia: Pedro predicó, luego fue el primero. Luego reprueba mi creencia sobre Matías y, sin embargo, no disuelve mi argumento. Matías y los demás apóstoles fueron obispos y no fueron ordenados por Cristo, ni se ordenaron a sí mismos; por tanto, no fueron ordenados sino por Pedro, a quien Cristo constituyó pastor universal al apacentar las ovejas. Por lo cual crea conmigo, o responda a ello.

En cuanto a la excusa de Pedro, me imputa cosas nunca pensadas, que yo habría separado el primado de las obras de administración de los demás apóstoles, cosa que nunca fue de mi intención; antes bien, le alegué al bienaventurado Gregorio: que los demás obispos han sido llamados a participar en la solicitud, no en la plenitud de la potestad.

Sobre Bernardo, acerca del tránsito de Pedro sobre el mar, dice que nada hace al primado. De lo cual me admiro, cómo puede decirlo, si lee a Bernardo. Pues Bernardo, por su intención, quiso probar a partir de esto el primado de Eugenio sobre los demás, y que todo el mundo debe estarle sujeto, no así a los demás obispos y apóstoles.

Sobre el estáter quiere sacar de la letra siguiente lo contrario, lo cual, sin embargo, no hizo, sino que Cristo reprimió la murmuración de los apóstoles. Pero yo no tomo las palabras de la sagrada Escritura por mi propia cabeza, \marginpar{\small[321]} sino según los santos padres, pues a partir de esto el bienaventurado Ambrosio quiso que se entendiera que Pedro habría de ser superior a los demás, lo cual también Agustín atestigua expresísimamente en las cuestiones sobre los Evangelios. Por lo cual no nos apoyemos en nuestro sentir, sino en los santos padres.

Sobre la indefectibilidad de la fe pasó de largo, aduciendo, sin embargo, aquello de Agustín de que la fe, que Pedro negó, el ladrón la conservó en la cruz. Y yo paso de largo sobre aquel dicho. Sepa, sin embargo, el reverendo padre que san Cipriano, apoyado en las palabras de Cristo, quiso que ninguna herejía surgiera en Roma como en las demás Iglesias, como ayer dije, que en la sede romana hubo doscientos

y doce papas y, sin embargo, antes de su remoción ninguno erró en la fe de modo definitivo y dictando sentencia.

En cuanto a Teófilo, a quien niega por su propio arbitrio sin aducir autoridad ni razón alguna, me admiro de que desprecie así a los padres eclesiásticos quien es tan hostil a los sofistas. Adujó a Agustín, quien no se opone a Teófilo en que Cristo habló del género de martirio. Pero enseñe dónde negó Agustín el orden del magisterio, cuando en el mismo tratado 124 sobre Pedro confiesa que él fue por gracia más abundante el uno y mismo primer apóstol. Y en la carta 162 dice: En la Iglesia romana, dice, siempre estuvo vigente el principado de la cátedra apostólica. Por lo cual no me pida que no defienda nuestra conclusión, que reputo verdadera.

Lo que trajo sobre Pedro enviado a Samaría, y por ello arrebatado el primado, como si no supiéramos que alguien puede enviarse a sí mismo, como el Hijo de Dios enviado por el Padre y por sí mismo, no ha empezado a ventilarse ahora por vez primera, sino que todo fue resuelto hace tiempo por Alfarus, Juan de Torquemada y otros. Y si Pedro cedió ante Santiago en Jerusalén en el concilio, por razón del lugar y de la ancianidad, no le arrebata el primado a Pedro. Pero, para parecer que hace algo, el reverendo padre atribuye a Pedro el primado de honor, acaso como a un ambicioso, aunque oyó de mí que, según Cipriano, Jerónimo y Anacleto, los apóstoles fueron de igual honor. Luego el reverendo padre atribuye a Pedro un primado de honor que los santos padres niegan, y en cambio niega a Pedro el primado de régimen que la concorde sentencia de los santos padres y el concilio de Constanza atribuye a Pedro.

Cese, pues, ruego, de gloriarse de estar apoyado en mayores autoridades contra tantos santos padres, contra tan célebre concilio, pues ya mostré muchas veces que las autoridades aducidas por él no niegan el primado según la sentencia de ningún santo padre. Por lo cual, más bien, con los griegos que al fin se hacen sabios, como se dice de los troyanos, aprenda por fin que el pontífice romano, verdadero vicario de Cristo, ocupa el primer lugar en el orbe, como los propios griegos, el emperador, el patriarca y los primados lo confesaron e hicieron obediencia a la Iglesia romana en el año 1439, el 22 de noviembre, y así se haga la paz en nuestros días. Todo lo cual, como arriba, encomiendo al juicio de aquellos a quienes toca y tocará. Solo a Dios la gloria. Eck.

\marginpar{\small[322]} El D. D. Martín.

Casi todo lo que ha dicho el egregio D. D. no me desagrada, en especial aquello de Agustín de que el principado de la cátedra apostólica estuvo siempre vigente

en la Iglesia romana, si hubiese añadido una palabra, a saber, el principado de potestad sobre todos los obispos, y creo que Bernardo retuerce el texto de Pedro caminando sobre el mar, pues la consecuencia atribuye manifiestamente aquella caminata a la fe, no al primado, al decir Cristo, cuando Pedro se hundía: Hombre de poca fe, ¿por qué dudaste [Mt 14 v. 31]? Y en cuanto a lo que desea, que yo crea con él que los demás apóstoles fueron ordenados obispos por el apóstol Pedro, de momento no me place obedecer, pues todos fueron obispos igualmente como Pedro, según prueba el texto aducido por Pedro acerca de Judas: Y su episcopado tómelo otro [Hch 1 v. 20; Sal 109 v. 8]. Esto sobre aquello.

\paragraph*{El D. D. Eck.}

Sobre el episcopado, en cuanto a que también Judas fue obispo según aquello del Salmo: Y su episcopado tómelo otro, digo que allí episcopado se usa por apostolado y no por el orden episcopal. Pues para la mayoría es dudoso si Judas llegó a ser siquiera sacerdote. Pero esto nada hace al caso. Sin embargo, Judas salió tras tomar el bocado; por lo cual, cuando en la cena Cristo ordenó sacerdotes a los discípulos, Judas no fue consagrado obispo. Lo encomiendo para que se juzgue, u otra cosa.

Fue renovada la disputa y enseguida se retomó la materia del purgatorio en torno a la novena posición.

El D. D. Eck comenzó como oponente.

No sólo entre los teólogos, sino también en la sagrada Escritura y en los santos padres se halla que esta vida presente es estado y término de merecer y desmerecer. Por lo cual los que están en el purgatorio ya no pueden merecer y así tampoco recibir mayor gracia. Que esto es así lo pruebo por Jer 25 [v. 14]: Les daré según sus obras y según los hechos de sus manos; así también dice Pablo: Es preciso que todos nosotros seamos manifestados ante el tribunal de Cristo, para que cada uno reciba lo propio del cuerpo, según hubiere obrado, sea bueno, sea malo [2 Cor 5 v. 10]. Conociendo, pues, el temor del Señor, persuadimos a los hombres.

Añado a Jerónimo, para descender por grados, entendiendo allí al apóstol: Esto dice, digo, que mientras establecidos en el mundo peregrinamos

lejos del Señor, con buena conducta procuremos agradarle en el futuro, no, como algunos piensan, que después de salir del cuerpo allí merezcamos algo obrando. Lo mismo confiesa la glosa ordinaria allí: En vano, pues, se promete el hombre después del cuerpo lo que descuidó conseguir aquí. Así también lo entendió Ambrosio: que sólo los hechos del cuerpo son premiados en otra parte.

De lo cual consta que, como la gracia aumenta según la correspondencia de los méritos y las almas en el purgatorio no podrán obrar nada meritoriamente por lo alegado, tampoco \marginpar{\small[323]} podrá aumentar en ellas la gracia. Lo suscribe clarísimamente el bienaventurado Agustín en el Enquiridión cap. 109 y 110, y repite las mismas palabras principalmente en la cuestión 2 de las ocho cuestiones a Dulcicio después de muchas cosas. Por lo cual, dice, aquí se adquiere todo mérito por el que alguien pueda después de esta vida ser aliviado o agravado. Y nadie espere merecer ante Dios, una vez muerto, lo que aquí descuidó. Cesa, pues, después de la muerte el aumento de la obra meritoria; cese también el aumento de la gracia correspondiente a los méritos.

\paragraph*{El D. D. Martín.}

El egregio D. D. opuso dos autoridades de la divina Escritura y cuatro de los santos padres. La primera, de Jeremías, donde Dios dice: Les daré según sus obras. Me admiro bastante de que el egregio D. D. juzgue que éstas y otras semejantes pugnan contra mí, cuando yo en mi resolutorio aduje muchas más en esta forma y expuse cómo nada pugnan contra mí o igualmente pugnan contra el purgatorio. Pues como la divina Escritura en su totalidad no contiene absolutamente nada sobre el purgatorio, sino que todo lo dice o del infierno o del cielo, quiero con esta única respuesta haber respondido a todas las autoridades, que no hacen al propósito, y en consecuencia tampoco todas las exposiciones de todos los padres, a no ser que mencionen expresamente el purgatorio. Pues es creíble que, cuando tratan las santas Escrituras, sigan también el sentido de la santa Escritura. Y por eso con otras razones y autoridades conviene que se enseñe que las almas están ciertas y que no aumenta en ellas la caridad.

\paragraph*{El D. D. Eck.}

Aunque el reverendo padre dice haber alegado más autoridades en el resolutorio y haberlas expuesto, ya las vimos. Vimos también que su glosa al texto nada aporta, lo cual ya mostraremos. Esto tampoco

aceptamos: que en la sagrada Escritura nada se halla sobre el purgatorio; dicho que ciertamente sería favorable a griegos y picardos, pero adverso a la fe cristiana, como rectamente recuerda en el resolutorio; sino opongamos su única solución al filo, por la que dice que lo aducido nada hace al propósito. ¡Con qué elegancia respondió! El Apóstol dice: Cada uno recibirá según lo que hubiere obrado en el cuerpo; el D. padre dice que el hombre recibe también según lo que hubiere obrado en el purgatorio. Agustín dice que todo mérito se adquiere aquí, y el D. padre dice en contra que también en el purgatorio se adquiere mérito. Jerónimo dice que, después de salir del cuerpo, nada merecemos ante Dios; el D. padre dice en contra que, después de salir del cuerpo, aún merecemos en el purgatorio; así de Ambrosio. Juzguen ya vuestras ilustrísimas dominaciones, excelencias y prestancias si lo aducido por mí nada hace al propósito, cuando directamente contraría lo dicho por el padre. Pido, pues, que realmente lo resuelva o muestre cómo no hace al pro-\marginpar{\small[324]}pósito, y no se admire de que también yo le voy a aducir más autoridades, para que sea manifiesto que esto no es de teologastros, sino de verdaderos teólogos.

\paragraph*{El D. D. Martín.}

Yo, que creo firmemente, más aún, osaría decir que sé que hay purgatorio, fácilmente me persuado de que en las Escrituras se hace mención de él, como aquello de Mt [12 v. 32] que aduce Gregorio en el Diálogo: No se le perdonará, ni en este siglo ni en el futuro, queriendo que ciertos pecados se perdonen en el purgatorio.

Admito también aquello de 2 Mac [12 v. 45 f.]: Santo y saludable es el pensamiento de rogar por los difuntos. Pero quiero esto: que en toda la Escritura no se halla memoria del purgatorio que pueda sostenerse en contienda y convencer, pues también el libro de los Macabeos, al no estar en el canon, es poderoso para los fieles, pero contra los pertinaces nada vale. Y el dicho de Gregorio facilísimamente se elude, que ni aquí ni en el futuro se perdone el pecado, esto es, nunca. Por eso no quiero que se me tenga por sospechoso de favorecer a bohemios y griegos. Esto por lo primero.

A la réplica del señor D. digo nada distinto de lo que ya antes dije. Pues repite lo mismo, y por eso también respondo lo mismo. Dije, pues, en

las sagradas letras nada se halla sobre el purgatorio, y por eso ni ellas ni sus exposiciones pueden aplicarse al negocio del purgatorio. Pues lo devuelto es o bien a los buenos o mal a los malos después de la muerte. En el medio están los que están en el purgatorio, de quienes no hablan dichas autoridades, a quienes ni se ha devuelto bien ni mal, y así el purgatorio pasa siempre por en medio de las autoridades.

\paragraph*{El D. D. Eck.}

Que se excuse diciendo que cree y sabe que hay purgatorio, lo acepto bien y ya lo sabía antes. Pero que niegue que el purgatorio se prueba en la sagrada Escritura, lo rechazo como falso y grecánico, permaneciendo en el propósito, aunque ni aquello ni otras cosas mucho mayores puedan convencer contra protervos y pertinaces, como constó en Arrio y otros. Pero que por eso se diga que no se halla en la Escritura, porque no pueda convencer al pertinaz, es del todo capcioso y adverso a nuestra religión.

Sobre los libros de los Macabeos, de los que dice que valen para los fieles pero que no están en el canon, también esto, digo, es falso, pues aunque entre los hebreos no estuvieron en el canon, la Iglesia, sin embargo, los recibió en el canon, como atestigua el padre del señor padre, Agustín, en el libro de la ciudad de Dios, y san Ifón inserta en sus decretos la constitución por la que la Iglesia recibió aquellos libros en el canon.

Que el reverendo padre quiera persistir en su propósito de que lo alegado por mí no hace al propósito, cuando, sin embargo, directamente le contraría. En primer lugar, lo que dice de que aquellas autoridades miran al cielo o al infierno no ayuda, porque directamente hacen mención del estado de merecer y desmerecer y definen que este \marginpar{\small[325]} estado termina con la muerte. En cambio, el reverendo padre, sin apoyarse en ningún santo doctor, osa prorrogar a las almas que salen del cuerpo el término de merecer y afirmar que los méritos del purgatorio son remunerados en el cielo, y la sentencia apostólica clama que sólo se remuneran las obras que el hombre obró en el cuerpo. Por lo cual examiné lo dicho en su resolución, pero no pude aceptarlo como insuficiente, como antes abrí aquellas contrariedades.

\paragraph*{El D. D. Martín.}

En primer lugar, en cuanto a lo que dice el egregio D. D. de que no por eso ha de negarse algo de la Escritura porque los pertinaces no puedan ser convencidos, lo dice óptima y verdaderísimamente, pero yo hablo de aquellos pertinaces que con nuestra propia

autoridad y con nuestro propio dardo pueden 'atravesarnos'. Pues es evidente que el libro de los Macabeos 'pertenece' al Antiguo 'Testamento'. Cuando, pues, 'san Jerónimo' hubiere descrito el canon hebreo y defina que sólo valen en contienda los libros que son del canon, y haya sido recibido en esta sentencia suya, fácilmente seremos perforados por nuestro propio dardo, a no ser que persuadamos a los fieles.

En segundo lugar prueba que el libro de los Macabeos ha sido recibido en el canon; contiende sobre una equivocación, y fácilmente nos concordaremos. Sé que la Iglesia recibe este libro, y esto lo dije, pero la Iglesia no puede atribuir más autoridad o firmeza al libro que la que por sí mismo tenga, así como aprueba y rechaza los opúsculos de los demás padres, pero no por eso confirma todo ni hace mejores las cosas que en ellos se dicen. Pase, pues, esa equivocación de canon y canon.

En tercer lugar dice que las autoridades aducidas hacen mención directa del estado de merecer y que éste termina con la muerte. Respondo como antes: Porque nada recuerdan del purgatorio, sino sólo del cielo o del infierno, por eso el que muere para el infierno no merece y el que muere para el cielo no merece, y así termina el estado de merecer de uno y otro; de otro modo, con esas mismas autoridades podría probarse evidentísimamente el purgatorio.

En cuarto lugar dice que yo no me apoyo en ningún doctor y que prorrogo el término de merecer, y que afirmo que los méritos del purgatorio son remunerados. Lo confieso, pues esto lo hice para disputar y oír cosas mejores que las que yo sepa. Pues yo nada sé del purgatorio, sino que las almas allí padecen y han de ser ayudadas con nuestras obras y oraciones, dispuesto a ser enseñado humildemente, si pueden transmitirse más cosas sobre él mismo. Por lo cual el Apóstol aducido contra mí, que sólo se remuneran las obras obradas en el cuerpo, ha sido tomado por el señor D. en relación al purgatorio, como ya dije que no puede tomarse, sino sólo al cielo o al infierno. Si se observa esa relación, es patente que no pugnan contra mí, o demostrarán el purgatorio.

Mañana a la hora séptima se continuará la disputa.

El sábado 9 del mes de julio a la hora séptima antes del mediodía.

\paragraph*{El D. D. Eck.}

Primero, cuando el D. padre dice que las autoridades aducidas no se entienden del purgatorio, place en cuanto a esto, porque ponen el mérito sólo en esta vida, luego excluyen el mérito \marginpar{\small[326]} en el purgatorio. Además es cierto que Agustín en el Enquiridión habla también del purgatorio, porque allí mismo prueba que las almas son aliviadas por los sufragios de los vivos. Y Jerónimo increpa a los que piensan que las almas despojadas merecen, el cual afirma el purgatorio, sobre el último capítulo de Isaías [66 v. 24].

Pero en cuanto a que dice que su conclusión es disputación, como no teniendo gran confianza en ella, me admiro muchísimo, cuando ya después de un año la ventiló e inscribió la disputa con magnífico título contra los errores nuevos y viejos, y afirma que la sentencia de los teólogos más recientes ni siquiera es verosímil para un necio y con desprecio los llama teologastros. Pero en cuanto a que imagina una equivocación del canon, no lo soporto, porque Agustín en aquel libro 18 de la ciudad de Dios no pudo equivocar el mismo término, máxime en el punto adversativo, al decir que no estuvo en el canon entre los hebreos, sino entre la Iglesia. Además es cosa explorada que, habiendo sido escritos más evangelios, por autoridad de la Iglesia cuatro fueron recibidos en el canon, y así el prólogo atestigua que los libros de los Macabeos fueron recibidos, pero que por la Iglesia son anotados entre las historias de los volúmenes divinos. — Pero mejor a la cosa.

Porque se funda en esto, que el purgatorio no está expreso en las sagradas letras, contra lo cual está el decreto del concilio florentino, que también los griegos asumieron abjurado el error, y después muchos lugares de la sagrada Escritura, como aquello: Pasamos por agua y fuego, y nos sacaste al refrigerio [Sal 66 v. 12]; aquello del Eclesiastés 4 [v. 14]: De la cárcel y de las cadenas sale para reinar; aquello de Mt 5 [v. 25]: Dijo el Señor: Ponte de acuerdo etc., no sea que seas echado en la cárcel. En verdad te digo, no saldrás de allí hasta que pagues el último cuadrante. Donde por cárcel entiende el lugar del purgatorio, del que el hombre no es liberado hasta haber satisfecho plenamente, como interpreta el bienaventurado Ambrosio en el lugar correspondiente sobre Lucas. Sin embargo, la razón de Gregorio, que retoma el bienaventurado Bernardo en el sermón 66 sobre el Cantar, es válida: que Cristo insinúa que algunos pecados se perdonan en el siglo futuro.

Pero el texto más abierto es 1 Cor 3 [v. 12—15]: Si la obra de alguno ardiere, sufrirá detrimento, pero él mismo será salvo, mas así como por fuego. Y el día del Señor lo declarará, porque será revelado en fuego, y el fuego probará cuál sea la obra de cada uno, y allí sobre la paja, el heno etc., donde el bienaventurado Ambrosio dice que el Apóstol habló expresamente del purgatorio. Lo mismo atestigua la glosa interlineal, y también la glosa ordinaria. Lo suscribe el bienaventurado Jerónimo en el libro 2 contra Joviniano, y también el bienaventurado Gregorio lo toma del purgatorio, lo que también hace el bienaventurado Bernardo sobre el Cantar. Y para que nadie me objete a Agustín en el Enquiridión, cap. acaso 67, quien dice que aquellas palabras pueden entenderse también del fuego del presente siglo, lea al mismo más completo allí mismo y en la cuestión 2 de las ocho cuestiones a Dulcicio, y hallará que Agustín toma aquellas palabras del Apóstol, como son manifiestísimas, también del purgatorio.

\marginpar{\small[327]} Por lo cual el purgatorio se afirma no sólo en los libros de los Macabeos, lo que ciertamente bastaría a la Iglesia, sino también en otras Escrituras, a no ser que Gregorio, Ambrosio, Agustín, Jerónimo y Bernardo sean del número de los teologastros; de otro modo no sé cómo el reverendo padre puede confesar que sabe que hay purgatorio, cuando quiere apoyarse tan válidamente en la letra de la sagrada Escritura, y cuando tan clarísimos son los testimonios de la Escritura en favor de los teólogos más nuevos. Él no podrá asignar ni un solo lugar en la Escritura donde se afirme que los méritos o la gracia aumentan en los que se purifican, cuando en cambio dice el Sabio, Ecl 11 [v. 3]: Y si el leño cayere hacia el austro o hacia el aquilón, en el lugar donde cayere, allí estará. Glosa ordinaria: Esto es, el lugar que aquí te hubieres preparado, entonces tendrás. Y porque en la casa de mi Padre, dice Cristo, muchas mansiones hay, al que muere se le asigna cierta mansión, más allá de la cual no puede ascender por sus méritos en el purgatorio. Así lo atestigua el Damasceno lib. 2 cap. 4: Conviene saber que lo que es en los hombres la muerte, eso es en los ángeles la caída, y lo declara en cuanto al término de merecer. Así el Sabio Eclesiastés 9 [v. 10]: Todo lo que puede tu mano obrar al instante etc., como para aquel sentido lo alegó aquel devoto y digno ministro de la Iglesia, Juan Capistrano. Así lo enseñó el Apóstol a Gál 6 [v. 7 ff.]: No erréis, Dios no es burlado. Pues lo que el hombre sembrare, eso también segará. Porque los que siembran en su carne, de la carne segarán corrupción etc., y haciendo el bien no desfallezcamos; luego, mientras tenemos tiempo, obremos el bien a todos. Glosa: lo que el hombre sembrare, a saber, con el trabajo de la vida presente. Luego no en el purgatorio

…se siembran los méritos, sino en la vida presente. Lo cual confirma también Crisóstomo por aquello de Jn 9 [v. 4]: «Me conviene obrar las obras del que me envió, mientras es de día; pues viene la noche, cuando nadie puede obrar». El día, dice Crisóstomo, es la vida presente. Por eso Agustín nos exhorta a que debemos obrar mientras vivamos, no sea que la noche nos sorprenda. Añádase como testigo también el bienaventurado Ambrosio, que entiende que el Apóstol habla del tiempo de la vida presente a nosotros concedido, para que vivamos justamente. Por lo cual dice el egregio salmista, Sal 104 [v. 23]: «Saldrá el hombre a su obra y a su labor hasta la tarde».

Por último, expresa es aquella sentencia de Agustín, además de las arriba alegadas, donde sin embargo hace mención del purgatorio; su sentencia es la del cap. 3 De fide ad Petrum: «El tiempo de alcanzar la vida eterna lo dio Dios a los hombres solo en esta vida». También Jerónimo, en el lib. 3 sobre aquello de Gal 6 [v. 5]: «Cada uno llevará su carga», muestra cómo los hombres son ayudados por sus méritos en esta vida, pero, llamados ante el tribunal, lo que sucede en la muerte, entonces ni Job ni Noé ni Daniel podrán interceder por nadie, sino que cada uno llevará su carga. Pues también bellamente enseña el bienaventurado Bernardo que en la muerte de cada hombre se realiza un juicio particular, lo cual no podría darse si [el alma] tuviera todavía plazo para acumular sus méritos en el purgatorio y aumentar la gracia; porque, como dice por el profeta Zac [4, v. 7], igualará gracia con gracia, devolverá gracia de gloria en proporción a la gracia de la fe. \marginpar{\small[328]} Y aunque aquellos testimonios no fuesen tan claros, siendo así que son clarísimos, esto solo debería apartar del propósito al reverendo padre: que un teólogo, sin fundamento teológico, a su propio arbitrio, no siendo juez ni asesor del juez, quiera prolongar a las almas que parten el plazo de merecer. Por lo cual, en cuanto a esta partícula, parece que debe concluirse que las almas en el purgatorio padecen bastante y son purificadas de sus delitos, pero no merecen mayor gracia.

Sobre lo segundo principal, porque niega que las almas en el purgatorio estén ciertas de su salvación, esto de nuevo juzgo que no está fundado en ninguna autoridad de la sagrada Escritura. En cambio, de que estén ciertas de su salvación, aduzco primero aquello de Ap 5 [v. 1 ss.]: «Vi en la diestra del que está sentado sobre el trono… Y nadie fue hallado digno de abrir el libro, ni en el cielo ni en la tierra ni debajo de la tierra». En el infierno indudablemente no hay dignidad alguna de abrir el libro; por lo cual habla del purgatorio, en el cual también hay a veces santos varones que resplandecen con milagros, como del santo Pascasio atestigua el bienaventurado Gregorio en el diálogo 4, y la historia del santo Severino, obispo de Colonia, de nación, sin embargo,

suevo, afirma lo mismo. Y más abajo: «Y todos cantaban, en el cielo, en la tierra y debajo de la tierra». Mas quienes están como desesperados de su salvación no cantan; por lo cual, por el hecho mismo de que el alma sagrada afirma que allí hay canto, afirma que las almas están ciertas de su salvación.

Además, en el canon de la misa, donde se ora por los difuntos en el purgatorio, decimos así: «Acuérdate, Señor, de tus siervos, que nos precedieron con el signo de la fe y duermen en el sueño de la paz», etc.; «a ellos y a todos los que descansan en Cristo», etc. Si, pues, según el sacratísimo canon de la misa las almas descansan en Cristo, ¿cómo pueden estar en tal turbación y horror como de desesperación, según declara el reverendo padre en su resolutorio de la concl. 14 a la 20? Pues no veo cómo en tal horror, temblor, turbación y como desesperación, cosas que todas dicen grandísima inquietud, las almas que han de purificarse puedan decirse que duermen en el sueño de la paz, siendo así que todo lo dicho se opone a la paz. En cambio, con razón serán juzgadas por los verdaderos cristianos que descansan en paz, mientras aguardan seguras el término de su purgación. Estas son en parte las cosas que me movieron y hoy me mueven a disentir del reverendo padre, dispuesto sin embargo a estar al juicio e información de otros.

\paragraph*{El D. D. Martín.}

Tres puntos por orden ha desarrollado contra mí copiosísimamente el egregio D. D.: primero, que el canon de la Escritura comprende los libros de los Macabeos; segundo, que el purgatorio se prueba también por otros lugares de la Escritura; tercero, ha intentado mostrar que las almas están ciertas de su salvación.

Respondo: sobre los dos primeros no hay ninguna controversia entre mí y su señoría. Por eso no era necesario acumular tantas autoridades para aquello \marginpar{\small[329]} que quizá afirmo con más constancia que él mismo, como quien he profesado saber que existe el purgatorio. Pero aquí se debate el punto capital: si puede probarse que las almas en el purgatorio merecen y que la gracia en ellas aumenta; pasemos, sin embargo, por orden.

Primero dijo que, con las autoridades ayer aducidas, establecí que el mérito está solo en esta vida, luego se excluye en la vida futura, lo que él entiende del purgatorio, y así adapta las autoridades al purgatorio. Yo, en cambio, concedo que se trata de la vida futura, del infierno o del cielo. Adujó también a Agustín en el Ench. hablando del purgatorio, y a Jerónimo increpando a quienes afirman que las almas separadas merecen. Una y otra cosa concedo. Me reprende también por haber dicho

que yo disputo esta proposición como no teniendo confianza en ella, yo que con magnífico título inscribí que disputaría contra los errores nuevos y viejos, y afirmé que la sentencia de los más recientes ni siquiera es verosímil para un necio, y llamé con desprecio «teologistas» a esos.

Digo, como ayer dije, que todavía nada sé «del negocio de las almas y del purgatorio», y a esto llamo error: que algunos se atreven a afirmar lo que ignoran, y las opiniones de los padres, como del bienaventurado Tomás, de Buenaventura y semejantes, que ellos mismos no afirmaron, ellos las proclaman como ciertos artículos de fe. A estos yo los llamo teologistas y no teólogos. Las opiniones deben tratarse en las escuelas; al pueblo deben predicarse las palabras y obras de Dios. Sal 19 [v. 2]: «Los cielos narran la gloria de Dios, y las obras de sus manos». No condeno yo, pues, las opiniones de los óptimos padres, sino que resisto a aquellos artífices que de las opiniones de los hombres nos forjan artículos de fe, lo cual no es oficio de buen teólogo.

Que yo haya usado equívocamente el canon contra Agustín, lib. 18, cap. 36, me obligaron a ello el divino Jerónimo y Eusebio, reseñando en la historia eclesiástica las autoridades de los antiguos. Por eso subsiste la equivocidad, en tanto que de un modo siente Agustín y de otro Jerónimo con los más antiguos acerca del canon, y por consiguiente no queda ninguna fuerza del argumento en la contienda. Si el prólogo de Jerónimo enumera los libros de los Macabeos entre los volúmenes divinos, no lo recuerdo; paso por alto aquello de que los cuatro evangelios han sido recibidos por autoridad de la Iglesia. Pues esta será otra materia.

Después de esto dice que es contra el concilio Florentino que el purgatorio no esté expreso en la Escritura. Respondo que el concilio no puede hacer que sea de la Escritura lo que por su naturaleza no es de la Escritura, como tampoco la Iglesia pudo hacer los evangelios, aunque aprobó los evangelios.

Por lo cual veamos las autoridades. Primero aquella de Sal 66 [v. 12]: «Pasamos por fuego y agua». Respondo: no vale para el purgatorio; habla de las persecuciones de los santos, como en otros muchos lugares: Sal 17 [v. 3]: «Me examinaste con fuego»; asimismo 26 [v. 2]: «Abraza mis riñones

y mi corazón», y 1 Pe 1 [1, v. 6 s.]: «Un poco, si ahora es preciso entristeceros en varias \marginpar{\small[330]} tentaciones, para que la prueba de vuestra fe sea mucho más preciosa que el oro, que por el fuego se prueba». Y en breve, aquel es un tropo vulgarísimo de la Escritura: entender por fuego y agua las tribulaciones. Por eso están demasiado apegados a letras y sílabas, como me imputa el egregio D. D., quienes toman el fuego por el purgatorio.

Asimismo también aquello del Eclesiastés: que alguien sale de la cárcel y de las cadenas al reino. Bastante claro es el texto que disputa sobre la vanidad de este mundo: que por giro del azar el que ahora es siervo se hace rey, y el que es rey se hace siervo. Con estas y semejantes autoridades, si lucháramos contra quienes [el purgatorio] niegan, no haríamos de nosotros y de la Iglesia sino ludibrio ante los adversarios; aunque yo gustosísimamente admita todas estas cosas, y cuantas haya semejantes.

Tal es también aquello de Mt 5 [v. 25]: «Ponte de acuerdo con tu adversario, no sea que te entregue al verdugo», donde dice que por la cárcel se entiende, según Ambrosio, el lugar del purgatorio; lo admito de buen grado. Pero como otros padres lo expusieron en sentido diverso, sobre todo Agustín, que lo expone del infierno del que nunca se saldrá, la autoridad se hace dudosa: persuade a los fieles, pero no convence a los que resisten; por callar que la ilación del texto ni siquiera consiente que se entienda del purgatorio, pues habla del que consiente con el adversario y del que disiente, queriendo que el que disiente peque condenablemente y contra el precepto de Cristo, cosas que pertenecen al infierno, no al purgatorio.

La razón de Gregorio, dice, consentida por Bernardo, es válida: que Dios perdona algunos pecados en el futuro. Respondo: es verdadera entre fieles, pero inválida ante los que resisten.

Después de esto dice que es texto clarísimo aquel de 1 Cor 3 [v. 15]: «él mismo se salvará, pero como por fuego», y que Ambrosio y Jerónimo, la glosa ordinaria y la interlineal, Gregorio y Bernardo entendieron al Apóstol del purgatorio, aunque Agustín diga en algún lugar que las mismas palabras pueden entenderse del fuego del presente siglo. Respondo: tan poco clarísimo es este texto, que hoy confieso, tras haber examinado mucho, que todavía no sé el sentido genuino de Pablo, tan variantes son los intérpretes; aunque de buen grado admita para mí el del purgatorio. Pero como el Apóstol dice abiertísimamente que la obra de cada uno será probada por el fuego, fuego en el que dice que será revelado el día del Señor, que declarará la obra de cada uno

la obra, como son palabras abiertas, a mi tenue juicio parece hablar del fuego de la conflagración y del juicio extremo, o, como Agustín lo tropologiza, del fuego de la persecución temporal, por la que sobre todo se prueba la fe y la doctrina de la fe, y cuantas cosas sobre estas se edifican.

Por lo cual todavía nada manifiesto se tiene de las sagradas letras sobre el purgatorio que valga en la contienda. No tengo, pues, a los beatísimos padres por teologistas, porque con ellos confieso el purgatorio; ni ellos sancionaron como artículos sus opiniones e ignorancias sobre el estado de las almas, como hacen los teologistas.

\marginpar{\small[331]} Se admira el D. D. de cómo puedo saber que existe el purgatorio, no teniendo yo ni un solo lugar en las Escrituras a mi favor, siendo así que estos, como dice, clarísimos testimonios de la Escritura juzga que favorecen a los teólogos neotéricos. Respondo: no es necesario confesar por qué vía haya conocido el purgatorio, o cualquier otra cosa. Además, estos clarísimos testimonios han sido aducidos a favor del purgatorio, no a favor del estado de las almas en el purgatorio.

Ya aquella autoridad de Ecl 11 [v. 3]: «En cualquier lugar donde cayere el leño, ya al austro ya al aquilón, allí estará», donde la glosa ordinaria entiende: el lugar que aquí te hayas preparado, entonces tendrás, no sé con qué ingenio se arrastra al propósito; pues si por lugar preparado y que se ha de tener entiende el purgatorio, habrá que permanecer eternamente en el purgatorio; pero si entiende que lo que aquí hayas preparado, esto es, merecido, entonces no hace a la autoridad de Ecl 11, que nada habla del mérito, sino de la muerte del hombre.

Paso por alto aquello de que Cristo dice que en la casa de su Padre hay muchas mansiones, y que al que muere se le asigna cierta mansión, más allá de la cual no puede ascender. Todo es a mi favor. Sé que a cada uno se le asigna una mansión después de la muerte, pero no llega enseguida a la mansión después de la muerte, a no ser que tome de nuevo la mansión por el purgatorio, y así excluya a los moribundos del reino eterno.

Asimismo también aquello del Damasceno: lo que para los ángeles la caída, eso es para los hombres la muerte. Respondo: si, pues, después de la muerte cayeren en el purgatorio, se sigue que eternamente serán purgados; o es preciso entender al Damasceno, como todo lo precedente, de las dos últimas mansiones, la condenación y la bienaventuranza, y no del purgatorio.

Así también aquello de Ecl 9 [v. 10]: «Todo lo que pudiere tu mano, óbralo al instante, porque ni razón ni sentido hay en los infiernos, adonde tú te apresuras». Si lo entiende del purgatorio, de nuevo el purgatorio será infierno; luego habla solo del infierno, sin memoria del purgatorio.

A Juan Capistrano lo admito de buen grado, pero fuera de la contienda.

Ya que lo que dice el Apóstol, Gal 6 [v. 7], habla contra nuestra sentencia: «Lo que el hombre sembrare en la carne, eso también segará», y aquello de Jn 9 [v. 4]: «Viene la noche, cuando nadie podrá obrar», lo que Ambrosio expone del tiempo de la vida presente, en el que merecemos, claro es que nada hace al caso. Pues el Apóstol no habla del purgatorio, sino del juicio extremo. Y Cristo, por la noche —sea dicha con reverencia salvo la de Crisóstomo— entiende la infidelidad, como es claro por el texto próximo: «mientras yo estoy en el mundo, luz soy del mundo», queriendo decir que fuera de la fe de Cristo nadie puede obrar bien; sin embargo admito a Crisóstomo, porque mira al juicio extremo.

La autoridad de Sal 104 [v. 23]: «Saldrá el hombre a su obra hasta la tarde», se permite arrastrarla en sentido figurativo a la vida y muerte del hombre, pero en sentido genuino, y que luche en la contienda, habla de la admirable disposición de Dios, \marginpar{\small[332]} que así dispuso al hombre: que por la mañana salga hasta la tarde del día natural a su obra. Y Agustín rectamente siente, en el cap. 3 De fide ad Petrum, que el tiempo de alcanzar la vida eterna les ha sido dado a los hombres solo en esta vida, porque, como dije, siempre miran a la vida futura, no al purgatorio.

Y Jerónimo, sobre Gal 6 [v. 5]: «Cada uno llevará su carga», bien muestra que los hombres son ayudados por sus méritos en esta vida, pero, llamados ante el tribunal, lo que sucede en la muerte, entonces ni Job ni Noé ni Daniel pueden interceder por nadie. Respondo que esta autoridad es robustísima contra el egregio D. D., porque si es verdad que el hombre en la muerte es llamado ante el tribunal, y entonces ni Job ni Noé ni Daniel pueden interceder, en vano ora la Iglesia por los muertos. Y así se niega el purgatorio. Por lo cual Jerónimo habla del tribunal novísimo, después de terminado el purgatorio.

A Bernardo también lo adujo, en cuanto que en la muerte del hombre se celebra un juicio particular, lo cual no se daría si [el alma] tuviera todavía plazo para merecer. Respondo: sea lo que fuere del juicio particular, la consecuencia no vale.

Paso por alto la autoridad de Zac [4, v. 7]: «Igualará gracia con gracia», porque se expone en sentido tropológico y bueno, pero no propio.

Al fin de este artículo dice: aunque estos testimonios no fuesen tan claros, esto, sin embargo, debería apartarme a mí, teólogo, a saber, que así sin fundamento, a mi arbitrio, no habiendo sido yo ni juez ni asesor, habría prorrogado a las almas que parten el término de merecer. Lo mismo devuelvo contra el D. D., no habiendo sido tampoco él ni juez ni asesor, y sin embargo define para las almas un estado cierto según sus opiniones, sin fundamento; sobre todo porque no quiere que esto sea solo opinión (lo que toleraría gustosísimamente), sino ciencia cierta.

A la hora segunda de la tarde continuará la disputa.

A la hora segunda de la tarde del mismo día.

El D. D. Martín continuó lo comenzado, como sigue.

Y como mi sentencia no parece bastante entendida por el egregio D. D., en cuanto pueda me expondré más abiertamente: que las Escrituras y los santos padres antiguos en sus sentencias miran a la vida futura, en la que o se salvan o se condenan las almas, sin pensar entretanto nada del purgatorio. Por eso aquellas muchas autoridades que niegan el estado de merecer después de esta vida no pertenecen al purgatorio. Que si aún no se entiende, fínjase que no hay purgatorio alguno, como tampoco estuvo su pensamiento en las mentes de ellos. Entonces todavía se sostendrán y serán verdaderas todas las autoridades aducidas: que antes de la muerte está la vida del merecer. Además, no hablo de los méritos de las almas en el purgatorio en el sentido de que obren algo, sino de que reciben gracia más amplia. Como entre todos consta \marginpar{\small[333]} que la culpa no se perdona, ni siquiera la venial, sin aumento de la gracia, y el divino Gregorio escribe que en el purgatorio se perdonan las culpas, en el diálogo 4, y se recita en la dist. 25, cap. Qualis, así, sin embargo, afirmo estas cosas de modo que las tengo por opinión,

más aún, para que lo confiese, por ignorancia. Creo que solo a Dios le es conocido el estado de las almas en el purgatorio.

Al tercer punto principal, sobre la ciencia de su salvación, donde adujo primero Ap 5 [v. 3], que nadie fue hallado ni en el cielo ni en la tierra ni debajo de la tierra digno de abrir el libro, queriendo entender por «debajo de la tierra» el purgatorio, como en el que «alguna vez» también estuvieron santos varones, como Pascasio y Severino:

Digo que esta glosa no tiene autoridad; por eso con la misma facilidad se desprecia con que se prueba. Confieso, ciertamente, que los santos estuvieron en el purgatorio; más aún, que nadie está en el purgatorio sino santo. Podría, pues, decir un porfiado que «debajo de la tierra» es lo mismo que el infierno o cualquier otra cosa, como que ni los demonios ni los hombres ni los ángeles pueden abrir el libro; igual que se dice que la triple máquina de las cosas honra a Dios, y según el Apóstol, tanto las celestes como las terrestres como las infernales doblan la rodilla; si es así que también los demonios tiemblan; más aún, como es el modo en el mismo libro, «debajo de la tierra» significa también a los muertos, como dice en otro cap. [20, v. 13]: «La tierra dio sus muertos, y el infierno dio sus muertos».

En cuanto a lo que seguidamente aduce: todos cantaban en el cielo y en la tierra y debajo de la tierra, y que los desesperados no pueden cantar, respondo: nunca dije que las almas desesperen, sino que, siendo una misma la sentencia de la Iglesia de que es igual la pena de los infiernos y del purgatorio, dije que son semejantes a los desesperados; del modo como se lee que también en esta vida algunos fueron tentados de desesperación, como en el Salterio en muchos lugares: «No apartes tu rostro de mí, y seré semejante a los que descienden al lago [27, v. 9; 28, v. 1]». Además, cantar al Señor no siempre es alegrarse y gozarse; más aún, el cántico nuevo es el cántico de la cruz, esto es, alabar y llevar a Dios en medio de las tribulaciones, y hasta en la muerte.

En cuanto a lo tercero, que adujo el canon de la misa, donde orando por los difuntos decimos: «Duermen en el sueño de la paz» y «descansan en Cristo», no ve, según dice, cómo puedan decirse que duermen en el sueño de la paz, si están en tal turbación, temblor e inquietud máxima. Por eso, interpretando este descanso de paz, dice que aguardan seguras el término de su purgación. Respondo: no me satisface esta glosa, y se prueba lo mismo por lo mismo. Pues que estén inquietos lo prueba lo que sigue en el canon:

«A ellos, Señor, y a todos los que descansan en Cristo, concédeles lugar de refrigerio, de luz y de paz», y aquello que todos oramos: «Dales, Señor, el descanso eterno». Asimismo: «Dales la paz», lo cual no puede entenderse como «dales aguardar seguras el término de su purgación», que es lo que el D. D. interpretó que es la paz. Más rectamente, pues, a mi juicio descansan en paz en cuanto al cuerpo. Pues dormir en las sagradas letras en paz significa descansar en el sepulcro. \marginpar{\small[334]} Y así es patente que el egregio D. D. no probó la ciencia de su salvación; más aún, en cierto modo quita las penas del purgatorio, mientras les atribuye paz.

\paragraph*{El D. D. Eck.}

Como nos ha sido fijado un tiempo estrecho para terminar hoy esta materia, de modo que por la calidad de la materia no puedo refutar las soluciones vanas ni tocar otros puntos del purgatorio, aduciré unas pocas cosas, para que fácilmente se reconozca que el reverendo padre ha evitado la médula de nuestra materia y ha buscado desvíos. De donde, primero, habiendo yo aducido a Agustín —que todo mérito aquí se adquiere, y que nadie espere alcanzar ante Dios, una vez muerto, lo que aquí descuidó—, y asimismo a Jerónimo, dice que una y otra cosa concede. Pero ¿qué obstinación es entonces esta, si concede con Agustín que todo mérito aquí se adquiere, para decir que la gracia en las almas, aun en el purgatorio, aumenta?

Ni vale en este caso su colirio para todas las autoridades, de las que dice que no hablan del purgatorio, pues el bienaventurado Agustín dijo aquello mismo tratando de las almas en el purgatorio, tanto en el Enchiridion como en el libro De cura pro mortuis agenda, c. 1. Pues después de la muerte no podemos hacer lo que aproveche, sino recibir lo que hicimos. En esta dificultad había caído Agustín: cómo podrían aprovechar a los muertos los sufragios, no mereciendo nada. Recurrió a esto: a que en vida merecieron lo que después de la muerte les aprovechase; por lo cual en vano se esforzaba en instruirme sobre el sentido, que bastante bien había yo comprendido. Ni su persuasioncilla tiene momento alguno ante un cristiano, porque las veniales no se perdonan sin aumento de la gracia. Pruébelo esto el reverendo padre. Con ninguna autoridad está esto cimentado; más aún, es falso: como el pecado venial no produce ofensa de Dios por misericordia de Dios, no es necesaria nueva gracia para su borradura, sino que basta que alguien padezca por él satisfacción.

En cuanto a lo que en segundo lugar se excusa del error, porque no veo aducida ni una letra por el D. padre a favor de su nueva doctrina: por eso con razón

debe reputarse sospechosa, siendo así que la sentencia de los teólogos más recientes se apoya en tantas autoridades de la sagrada Escritura y de los santos padres. Ni le excusa esto: que a algunos imputa hacer de las opiniones de Tomás o de Escoto como artículos de fe, cosa que a mí no me consta. Pues yo no aduje a Tomás o a Escoto, sino a las clarísimas lumbreras de la Iglesia junto con la sagrada Escritura; vea él mismo a qué teologistas acusa.

Tercero, a Agustín opone, contra los libros de los Macabeos, con más fuerza al bienaventurado Jerónimo; pero Jerónimo en ninguna parte niega que los libros de los Macabeos estén en la Iglesia dentro del canon de la Biblia; es más, en el prólogo consta constantemente haber sido anotado entre las historias de los volúmenes divinos. Por lo cual no debió equipararse a los opúsculos de los santos padres en el can. Sancta Rom. 15 dist.

Y habiéndole yo opuesto lo del concilio Florentino, respondió que el concilio no puede \marginpar{\small[335]} hacer que algo sea de la Escritura si no lo es. Esto, ciertamente, es verdad; pero ¿qué es esto: mancillar con tanta temeridad un concilio tan laudable, hasta decretar este absurdo? Siendo así que fueron varones doctísimos los de aquel concilio, prefiero creer al concilio, que es regido por el Espíritu Santo, que al señor Lutero; no porque el concilio haga que algo sea de la Escritura si no lo es, sino porque creo que el concilio tiene el sentido y la inteligencia de las Escrituras al decretar que esto es de la Escritura, lo cual en la Escritura se halla.

Vanida es su excusa de que siempre puede hallarse una exposición tal que el texto alegado no valga en la contienda contra los pertinaces, pues ese sería el resbaladero de todos los herejes, que siempre podrían aportar algún afeite de exposición con el que contendieran que las verdades católicas no están expresas en la sagrada Escritura. Así hasta hoy duraría aún la pérfida herejía de Arrio: que el homousía no podría probarse tan expresamente por las sagradas letras que en la contienda no escapara un pertinaz con cualquier afeite.

Igualmente bien sabemos que la Iglesia no puede hacer los evangelios; sin embargo, la Iglesia hace que, dejados los evangelios de Nicodemo, de Bartolomé, de Tomás y de otros, solo a cuatro prestemos fe indudable; en cuyo juicio de la Iglesia hay que estar en la elección de los evangelios. Así también en la inteligencia y exposición de las sagradas Escrituras.

Además, asumiendo las autoridades por mí aducidas desde el principio, quiso que fuera innecesario lo hecho, porque también él creería que el purgatorio existe,

no recuerda que, frente a tantas autoridades que fijan en la vida presente el término de merecer, quiso zafarse porque el purgatorio no se probaría a partir de las Sagradas Escrituras; dicho lo cual, con razón debí yo refutarlo, no fuera que los picardos y otros cismáticos lo tomasen a su favor, ellos que, como el venerable padre, no saben que el purgatorio existe. Dejadas a un lado algunas, contemplemos al menos dos citas.

En Mt 5 [v. 25] no quiso entender por cárcel el purgatorio, porque también Agustín, lo cual en absoluto se me ocultaba, entiende por cárcel el infierno. Añado, más allá del venerable padre, que Crisóstomo entiende la vida presente, pero que es más recta la sentencia de Ambrosio, o al menos no despreciable, él que por cárcel entiende el purgatorio. Las palabras de Cristo esto indican: que aquél debe pagar hasta el último cuadrante; pero en el infierno no hay pago alguno, como tampoco hay redención. Ni el subterfugio añadido puede rechazar la exposición de Ambrosio, a saber, que el que ha de ser arrojado a la cárcel pecó mortalmente, puesto que también por los pecados veniales y por los pecados mortales, aunque ya contritos, se sufre pena en el purgatorio, lo cual recibimos de los Macabeos, cuando dice: Santa y saludable es, pues, la resolución de orar por los muertos, para que sean liberados de los pecados. Pues estos que habían sido muertos, y por quienes Judas Macabeo hizo ofrendas, habían pecado mortalmente a causa de los despojos de los ídolos, aunque se cree que se arrepintieron en \marginpar{\small[336]} la misma matanza, según la glosa ordinaria allí mismo, y aquello del Sal [78. v. 34]: Cuando los hacía morir, lo buscaban.

En cuarto lugar, cuando hube aducido el clarísimo lugar del apóstol Pablo, 1 Cor 3 [v. 12ss.], que Agustín, Ambrosio, Jerónimo, Gregorio, Bernardo, Isidoro con la glosa ordinaria entienden del fuego del purgatorio, confiesa el reverendo padre que no tiene el sentido genuino de este lugar; por eso, despreciados tantos padres, aportó una sentencia nueva: que el apóstol habla del fuego de la conflagración, porque el apóstol hizo mención del día del Señor, y de que el fuego probará la obra de cada uno; yo, que siempre oí que debe creerse sobre todo a los santos más antiguos, abrazo con veneración la sentencia de los santos padres y no acepto aquella glosa nueva ni las palabras del D. Lutero, salvo que lo pruebe con autoridad de la Sagrada Escritura. Ni le ayudan las consideraciones que pondera en esta ocasión, que también los santos padres leyeron bien. Pues aunque peculiarmente el día del juicio final se diga día del Señor, como también recuerda Bernardo, sin embargo en la muerte de cada hombre, cuando se dicta el juicio, puede decirse día del Señor, y no acepta una distribución tan exacta y capciosa, como suelen hacer los sofistas, en la palabrita «de cada uno». Pero sobre esto, reciba a quienes sobreedificaron paja, madera y heno sobre el fundamento, como también Agustín complica la distribución de Jn 1 [v. 9]: Ilumina a todo hombre que viene a este mundo. De donde es invicta la sentencia apostólica: que quienes sobre el fundamento edifican madera, heno, paja, se salvan, pero por el fuego de la purificación, de modo que ningún mal quede impune.

A aquello del Eclesiastés sobre la caída del leño responde: Si habla del purgatorio, así se probaría que el alma permanece perpetuamente en el purgatorio. Digo que sintió óptimamente la glosa: que por leño se entiende el hombre, como también en Ez 31 [v. 8] significa la criatura racional: Todo leño del paraíso etc., y por caída la muerte; y sin embargo no se sigue que por eso permanezca perpetuamente en el purgatorio, sino que, como expone egregiamente el bienaventurado Jerónimo, por el austro se significan el bien, y por el aquilón el mal y la condenación. Por eso el sabio no entendió aquellos lugares, sino el estado: que en la caída o es bueno, y así perseverará sin aumento de gracia por la supresión de los pecados, o si es malo, perdurará en la malicia.

Sobre la mansión en Jn [14 v. 2] dijo que, inmediatamente después de la muerte, se asigna al alma una mansión cierta. ¿Cómo, pues, será cierta, si se produce un aumento de gracia mayor? Pues entonces necesariamente sobrevendrá una mansión más alta; también sería consecuencia que el que descendiera al purgatorio con muchos veniales estaría mejor que el que descendiera con pocos o el que muriera sin ninguno, porque se le haría mayor aumento de gracia para la supresión de los muchos veniales. Añado que sería nocivo y condenable orar por los muertos, y más les aprovecharía si perdurasen en las penas del purgatorio por el aumento de gracia mayor, como en algún lugar escribió el reverendo padre: que el alma liberada por sufragios ajenos es menos beatificada que si por sí misma padeciese lo suficiente en el purgatorio. Lo cual es no sólo contra la Sagrada \marginpar{\small[337]} Escritura: Santa, pues, y saludable es la resolución etc., sino también contra la observancia y la piedad de toda la Iglesia, que debemos tributar a los muertos, como el bienaventurado Agustín expone bellísimamente en el libro sobre el cuidado que ha de tenerse de los muertos y en el libro Enquiridión.

A lo de Damasceno respondió además que no habla del purgatorio, puesto que así perdurarían perpetuamente en el purgatorio. Y yo siento que ni Damasceno ni las demás autoridades hablan del purgatorio, pues de otro modo me dañarían y aprovecharían al venerable padre. Pero como hablan de la vida presente, en la que ponen el término de merecer y desmerecer, no puede extender el tiempo de merecer más allá del término establecido, esto es, la muerte. Por eso no acepto otra exposición que la que el propio autor quiera, el cual habla en el lugar alegado del término de merecer.

Así también puedo decir de los demás, y no puedo ser increpado con semejante invectiva como el reverendo padre, porque con la autoridad de los santos padres y de la Sagrada Escritura cierro el curso y el término de merecer con la muerte. Él, en cambio, sin apoyo alguno de la autoridad de la Sagrada Escritura, les extiende y prorroga el término, contra Agustín y contra Jerónimo, aunque sea propio de su modestia decir que esta sentencia suya es su opinión; pero entonces no debió llamar error a lo opuesto.

En cuanto a la autoridad de Jerónimo, intentó volverla contra mí: que ni Noé ni Job ni Daniel interceden por nadie después de la muerte, porque allí se habla del juicio final, lo cual no se prueba por la letra. Pues no sólo en el juicio final, sino también inmediatamente después de la muerte cada uno llevará su carga.

Finalmente, pasando a otra partícula, respondió a aquello de Ap [5 v. 3ss.] que por debajo de la tierra se entiende el infierno. Pero entonces lo siguiente sobre el cántico no puede sostenerse, porque los demonios y los condenados no cantan, sino que blasfeman y aúllan. Por eso por debajo de la tierra necesariamente puede entenderse el purgatorio y no los infiernos. Aunque no niego que en la Sagrada Escritura también se haga mención de los demonios en los infiernos, como en el apóstol: de los celestes, terrestres e infernales. Pero esto nada tiene que ver con el asunto.

Aportó además una glosa sobre aquella palabra «canta», según la cual sería llevar la cruz y alabar a Dios en las tribulaciones; glosa que, como forjada por el venerable padre en este lugar, desprecio con la misma facilidad. Más aún, Juan no tolera ser expuesto así, él que dice: Cantaban en el cielo, no en la tierra; pero en el cielo no llevan la cruz ni alaban en las tribulaciones, porque enjugó Dios toda lágrima de los ojos de los santos y ya no hay trabajo etc..

Pero aun si se concediera su exposición, como si fuera verdadera, aunque repugna mucho a la letra, todavía no quebrantaría este propósito, pues los que alaban a Dios en las tribulaciones después de la muerte tendrían certísimo signo de la futura bienaventuranza, mientras que los condenados, por el contrario, con odio consumado de Dios prorrumpirán en perpetua blasfemia del Creador. \marginpar{\small[338]} Por eso las almas que han de purificarse, seguras de su salvación, cantan y alaban a Dios, aguardando el término de la purificación.

Pero cuando hube aducido el canon de la misa, en el que se afirma que ellos duermen en el sueño de la paz, no tolera mi inferencia: si en el sueño de la paz, luego están seguros de su salvación y no en horror, temblor, pavor y casi desesperación. Aportó una doble razón: primero, porque sigue que les concedas un lugar de refrigerio, de luz y de paz; segundo, porque frecuentemente oramos: Dales, Señor, el descanso eterno, y: Dales la paz. Por eso dio otra glosa: que duermen en el sueño de la paz, esto es, en cuanto al cuerpo. No me contenta ni la objeción ni su nueva glosa, que no tiene apariencia alguna ni fundamento alguno. Pues dice el canon: Acuérdate de ellos etc., que duermen en el sueño de la paz, y más abajo: Y de todos los que descansan en Cristo. No dice: De los que descansan en el sepulcro. Además, ¿qué importa al alma que el cuerpo tenga paz o no, ya sea arrojado al mar o lacerado por las fieras, como atestigua el bienaventurado Agustín? ¿Qué hombre soporta que el cuerpo exánime y exangüe duerma en paz? De donde no admite de buen grado las distinciones en los teólogos más recientes, y sin embargo aquí distingue egregiamente entre la paz del cuerpo y la del alma.

Por eso ha de abrazarse la sentencia más verdadera: que las almas que han de purificarse duermen en paz y descansan en Cristo. De donde aquel descanso y aquella paz excluyen la turbación del alma, el horror y la casi desesperación, pues estarían inquietísimas las almas que fueran agitadas por esta turbación de la mente, hasta casi desesperar.

A lo aducido, en cambio, respondemos sin mucho gasto: Que oramos y suplicamos a Dios que les dé la paz, no la paz que ya tienen ahora, esto es, la seguridad de su salvación, sino la paz eterna, puesto que la paz es, como atestigua Bernardo, tanto respecto de la culpa como de la miseria. Dé, pues, Dios a las almas en el purgatorio la paz frente a las penas, como les confirió la gracia contra el pecado.

Por eso sigue siendo verdadera e inconmovible la verdad: que las almas en el purgatorio no casi desesperan de su salvación.

\paragraph*{El D. D. Martín.}

Contra las soluciones replica el egregio D. D. desde Agustín, que dice: Todo mérito se adquiere aquí, y tratando del purgatorio, que después de la muerte sólo recibamos lo que hicimos. Respondo brevísimamente: Luego o no hay purgatorio, o el purgatorio es lo que se recibe por la vida pasada. Por eso es patente que Agustín no puede entenderse del purgatorio. Pues no hicieron esto en la vida pasada, para recibir el purgatorio. Aunque también esta autoridad me patrocina, en que el mérito por el que merecen ser ayudadas lo merecieron aquí. Yo, en cambio, confieso que no puedo entender cómo son ayudadas las almas en el purgatorio y liberadas sin don alguno de gracia, por la sola quitación de las penas: entiéndalo quien pueda.

Luego dice que es falso que los veniales no se remitan sin aumento de gracia.

Cuando, pues, quiere el D. D. que los veniales puedan remitirse sin aumento de gracia, no lo concedo hasta que lo pruebe. Más aún, es contra el texto expreso, en Rom 7 [v. 14ss.], Gal 5 [v. 17], donde el apóstol dice: Vendido estoy bajo el pecado, ¿quién me librará de este cuerpo de muerte? Responde: La gracia por Jesucristo \marginpar{\small[339]}. Es cierto, en cambio, que el apóstol no estaba en pecado mortal. Lo demás lo dejo a la futura disputa entre los egregios señores doctores Andrés y Juan Eck. Y es erróneo esto: que el pecado venial no ofenda a Dios, pues a Dios desagrada toda inmundicia y reprueba al que quebrantare los mínimos, Mt 5 [v. 18s.]. Ni probó el D. D. que baste con que padezcan lo suficiente.

En segundo lugar me acusa de no haber aducido ni una letra en favor de mi sentencia, cuando él en favor de la sentencia de los teólogos más recientes adujo tantas autoridades de la Escritura y de los padres. Respondo: Tanto más sospechosa me es su sentencia cuanto más adujo, porque los trajo por los cabellos y retorciendo el cuello, como bastante dije en lo precedente. Menos peca quien duda en sus pensamientos que quien intenta establecer sus dudas con palabras divinas.

Dice también que no le consta quiénes establecieron como artículos de fe las opiniones de Tomás o de otros, porque adujo lumbreras de la Iglesia junto con la Escritura. Advierte, pues, que mire a quiénes llamo teologistas. Dije hoy y de nuevo digo: Permito que las opiniones sean opiniones y a nadie imputo esto como vicio, pero arrastrar la Escritura que se resiste al sentido ajeno para afianzar opiniones y permanecer pertinazmente en esto, esto digo que es oficio de teologistas.

